%%
%% AgentTelemetry: A Grounded Taxonomy and Empirical Evaluation
%% of Observability Primitives for Autonomous AI Agent Systems
%%
%% Target: MLSys 2027 / ICSE 2027 (10 pages + references)
%%

\documentclass[sigplan,10pt,screen]{acmart}

%% ---------- packages ----------
\usepackage{listings}
\usepackage{xcolor}
\usepackage{booktabs}
\usepackage{enumitem}
\usepackage{url}
\usepackage{multirow}
\let\Bbbk\relax
\usepackage{amssymb}
\usepackage{tikz}
\usetikzlibrary{positioning,arrows.meta,fit,backgrounds}

%% ---------- listings style ----------
\lstset{
  language=Python,
  basicstyle=\ttfamily\scriptsize,
  keywordstyle=\color{blue!70!black}\bfseries,
  commentstyle=\color{green!50!black},
  stringstyle=\color{red!60!black},
  showstringspaces=false,
  breaklines=true,
  breakatwhitespace=true,
  frame=single,
  framesep=2pt,
  xleftmargin=4pt,
  xrightmargin=2pt,
  numbers=left,
  numberstyle=\tiny\color{gray},
  numbersep=4pt,
  aboveskip=6pt,
  belowskip=4pt,
  columns=flexible,
}

%% ---------- metadata ----------
\title{AgentTelemetry: A Grounded Taxonomy and Empirical Evaluation\\of Observability Primitives for Autonomous AI Agent Systems}

\author{Krishna Chaitanya Balusu}
\affiliation{%
  \institution{Independent Researcher}
  \country{United States}
}
\email{krishnabkc15@gmail.com}

\begin{abstract}
Autonomous AI agents built on large language models are transitioning
from research prototypes to production deployments, yet the
observability infrastructure needed to diagnose their failures remains
fundamentally inadequate.  OpenTelemetry's GenAI semantic conventions
now cover LLM invocation, tool execution, retrieval, and agent
lifecycle, but leave five critical agent orchestration
phases---planning, reasoning, safety-check monitoring, inter-agent delegation,
and memory management---without any span-level representation.

We present \textsc{AgentTelemetry}, an open-source, OTel-native
observability library that introduces a \emph{grounded taxonomy} of
nine agent-specific span kinds---\texttt{AGENT}, \texttt{LLM\_CALL},
\texttt{TOOL\_CALL}, \texttt{PLANNING}, \texttt{REASONING},
\texttt{RETRIEVAL}, \texttt{GUARD\_RAIL}, \texttt{DELEGATION}, and
\texttt{MEMORY}---derived from systematic analysis of seven production
agent frameworks.

We make four contributions: (1)~the taxonomy itself, empirically
validated through an ablation study showing that removing any single
span kind makes at least one fault type undetectable; (2)~a
privacy-preserving multi-agent context propagation scheme extending W3C
Trace Context; (3)~a case study on 112~SWE-bench Lite instances (37\%
of the benchmark, across all 12~repositories) revealing that reasoning
loops account for 75\% of agent failures (95\% CI: [66\%, 82\%])---a
failure mode invisible to vanilla OTel but precisely characterized by
\textsc{AgentTelemetry}'s novel REASONING span kind---and a
telemetry-guided intervention that detects and breaks these loops,
improving the patch rate by $+11$~pp over a no-intervention control
(3-seed mean, $p{<}0.05$); and (4)~a pip-installable library
(3{,}700+ LOC, 78~tests) with adapters for seven frameworks,
including an end-to-end validated LangChain integration producing
correct spans from real \texttt{AgentExecutor} traces.
As a structural completeness check, a controlled fault-injection study
across 14~fault types, 6~LLM models, 7~frameworks, and 5~observability
conditions confirms that the span taxonomy is \emph{sufficient}:
\textsc{AgentTelemetry} achieves a Fault Detection Rate (FDR) of~1.000
under controlled injection---an upper bound demonstrating that every
fault type has a corresponding detection signal---compared to~0.429 for
vanilla OpenTelemetry and~0.000 without telemetry.
Instrumentation overhead is 11.7\,\textmu s median
(p99~$<$~30\,\textmu s) per span, negligible relative to LLM API
latencies of 500\,ms--10\,s.
\end{abstract}

\begin{CCSXML}
<ccs2012>
<concept>
<concept_id>10011007.10011074.10011099</concept_id>
<concept_desc>Software and its engineering~Software verification and validation</concept_desc>
<concept_significance>500</concept_significance>
</concept>
<concept>
<concept_id>10011007.10011074.10011111</concept_id>
<concept_desc>Software and its engineering~Software post-development issues</concept_desc>
<concept_significance>300</concept_significance>
</concept>
</ccs2012>
\end{CCSXML}

\ccsdesc[500]{Software and its engineering~Software verification and validation}
\ccsdesc[300]{Software and its engineering~Software post-development issues}

\keywords{AI agents, observability, distributed tracing, OpenTelemetry, LLM monitoring, fault detection, multi-agent systems}

\begin{document}

\maketitle

\noindent\textbf{Code and data:} \url{https://github.com/agenttelemetry/agenttelemetry}

%% ====================================================================
%% 1. INTRODUCTION
%% ====================================================================
\section{Introduction}
\label{sec:introduction}

Large language model (LLM) based autonomous agents represent a paradigm
shift in software architecture.  Unlike conventional services that
execute deterministic logic, an LLM agent pursues a goal through
iterative reasoning, planning, tool use, and
self-reflection---frequently delegating sub-tasks to other
agents~\cite{agentbench,agentboard}.  Frameworks such as LangChain,
CrewAI, AutoGen, and LlamaIndex have accelerated adoption, and
organizations are deploying agents for customer support, code
generation, and autonomous operations~\cite{aiops_survey}.

This proliferation has exposed a critical gap: \emph{there is no
standardized observability infrastructure for AI agent systems.}
Production teams cannot reliably answer fundamental diagnostic
questions:
\begin{itemize}[leftmargin=*,nosep]
\item \textbf{Why did the agent choose this tool?}  Decision
  attribution requires tracing from a tool invocation back through the
  LLM call that selected it and the reasoning chain that informed it.
\item \textbf{Where did the hallucination originate?}  Hallucination
  tracing requires comparing LLM output against retrieved documents,
  which demands semantic distinction between retrieval and generation
  spans.
\item \textbf{How much did this task cost across all agents?}  Cost
  attribution requires agent identity propagation combined with
  per-call token and pricing metadata.
\end{itemize}

OpenTelemetry (OTel)~\cite{otel} has become the standard for
distributed tracing in microservice architectures.  The OTel GenAI
semantic conventions define operation types for LLM calls
(\texttt{chat}, \texttt{text\_completion}), retrieval
(\texttt{retrieval}, \texttt{embeddings}), tool execution
(\texttt{execute\_tool}), and agent lifecycle
(\texttt{create\_agent}, \texttt{invoke\_agent}).  However, these
conventions address individual operations, not the higher-order
\emph{orchestration phases} that characterize autonomous agent
execution.  Five such phases---planning, reasoning, safety-check monitoring,
inter-agent delegation, and memory management---have no representation
in OTel GenAI, leaving them as opaque \texttt{INTERNAL} spans
indistinguishable from any other application logic.  This semantic
opacity makes automated fault detection substantially more difficult.

We present \textsc{AgentTelemetry}, an open-source, OTel-native
library that fills this gap with four contributions:

\begin{enumerate}[leftmargin=*,nosep]
\item \textbf{A grounded taxonomy of nine agent-specific span kinds},
  derived from systematic analysis of seven production frameworks and
  validated through an ablation study proving each kind is necessary
  for detecting at least one fault type (\S\ref{sec:taxonomy}).

\item \textbf{A lightweight multi-agent context propagation scheme} extending W3C
  Trace Context with agent identity baggage and three-level privacy
  controls (\S\ref{sec:propagation}).

\item \textbf{A SWE-bench case study demonstrating actionable observability:}
  on 112~instances, reasoning loops account for 75\% of agent
  failures---a failure mode invisible to vanilla OTel---and a
  telemetry-guided intervention improves the patch rate by $+11$~pp
  (3-seed mean, $p{<}0.05$), closing the loop from observability
  through diagnosis to remediation (\S\ref{sec:rq6}).

\item \textbf{Structural completeness via controlled fault injection:}
  across 14~fault types, 6~models, 7~frameworks, and 5~observability
  conditions, the span taxonomy achieves FDR\,=\,1.000---an upper
  bound confirming that every fault type has a detection signal---versus
  0.429 for vanilla OTel and 0.000 without telemetry (\S\ref{sec:evaluation}).

\item \textbf{An open-source, pip-installable library} (3{,}700+ LOC,
  78~tests) with seven framework adapters---one reference
  implementation providing full fault coverage, plus six
  adapter-validated stubs---using three distinct
  instrumentation strategies (\S\ref{sec:implementation}).
\end{enumerate}

%% ====================================================================
%% 2. BACKGROUND AND MOTIVATION
%% ====================================================================
\section{Background and Motivation}
\label{sec:background}

\subsection{Agent Execution Model}

An LLM agent's execution produces a deeply nested, non-deterministic
trace of cognitive operations.  A multi-agent research
assistant---comprising planner, researcher, and writer agents---may
produce dozens of spans across three or more agents, accumulating
thousands of tokens and non-trivial cost.  Debugging a failure requires
tracing backward across agent boundaries; without unified tracing,
operators resort to ad-hoc logging that is inconsistent across
frameworks and scales poorly.

\subsection{The OpenTelemetry Gap}

OpenTelemetry provides trace and span identifiers, context propagation,
and a rich exporter ecosystem.  However, its semantic layer was designed
for network-centric microservices.  The OTel GenAI semantic
conventions~\cite{otel} (all currently at Development stability) have
made significant progress: they define eight \texttt{gen\_ai.operation.name}
values---\texttt{chat}, \texttt{text\_completion}, \texttt{embeddings},
\texttt{retrieval}, \texttt{execute\_tool}, \texttt{generate\_content},
\texttt{create\_agent}, and \texttt{invoke\_agent}---covering
LLM invocation, vector retrieval, tool execution, and agent lifecycle.

Despite this breadth, the GenAI conventions address only \emph{individual
operation types} and leave a critical gap in \emph{agent orchestration
phases}.  Table~\ref{tab:gap} summarizes: five of nine
\textsc{AgentTelemetry} span kinds---\texttt{PLANNING},
\texttt{REASONING}, \texttt{GUARD\_RAIL}, \texttt{DELEGATION}, and
\texttt{MEMORY}---have no OTel GenAI equivalent at all.  A sixth
(\texttt{AGENT}) extends OTel's \texttt{invoke\_agent} with framework,
role, and orchestration semantics.  Without these distinctions, a
planning step and a summarization call are both opaque \texttt{chat}
spans; a safety-check outcome and a reasoning chain are both undifferentiated
\texttt{INTERNAL} spans.

\begin{table}[t]
\centering
\small
\caption{Agent concepts mapped to observability primitives.
$\bullet$\,=\,overlapping with OTel GenAI semantic conventions;
$\triangleright$\,=\,extends OTel GenAI; $\star$\,=\,novel, no OTel GenAI equivalent.}
\label{tab:gap}
\begin{tabular}{@{}llll@{}}
\toprule
\textbf{Agent Concept} & \textbf{OTel GenAI} & \textbf{AgentTelemetry} & \\
\midrule
Agent lifecycle    & create/invoke\_agent  & \textbf{AGENT} & $\triangleright$ \\
LLM invocation     & chat / text\_compl.   & \textbf{LLM\_CALL} & $\bullet$ \\
Tool execution     & execute\_tool         & \textbf{TOOL\_CALL} & $\bullet$ \\
Doc.\ retrieval    & retrieval, embeddings & \textbf{RETRIEVAL} & $\bullet$ \\
Task decomposition & (none)                & \textbf{PLANNING} & $\star$ \\
Chain-of-thought   & (none)                & \textbf{REASONING} & $\star$ \\
Safety-check monitoring & (none)               & \textbf{GUARD\_RAIL} & $\star$ \\
Inter-agent handoff & (none)               & \textbf{DELEGATION} & $\star$ \\
Memory read/write  & (none)                & \textbf{MEMORY} & $\star$ \\
Token count        & gen\_ai.usage.*       & \texttt{llm.input/output\_tokens} & $\bullet$ \\
Cost estimation    & (none)                & \texttt{llm.cost} & $\star$ \\
Agent identity     & gen\_ai.agent.*       & \texttt{agent.name, .role} & $\triangleright$ \\
\bottomrule
\end{tabular}
\end{table}

%% ====================================================================
%% 3. SPAN TAXONOMY
%% ====================================================================
\section{A Grounded Taxonomy of Agent Span Kinds}
\label{sec:taxonomy}

\subsection{Derivation Methodology}
\label{sec:derivation}

The nine span kinds were derived through a structured qualitative
analysis loosely modeled on open coding~\cite{strauss_corbin}: we
examined framework APIs, generated candidate span kinds from observed
execution phases, iteratively merged and refined candidates across
frameworks, and stopped when no new kinds emerged.  The process
proceeded in four stages.

\textbf{Stage~1: API inventory.}  For each of the seven frameworks, we
catalogued every public API entry point that represents a distinct
execution phase in an agent loop.  Table~\ref{tab:traceability} lists
the 25~entry points examined.  Callback-based frameworks (LangChain,
CrewAI, LlamaIndex) expose named hooks
(e.g., \texttt{on\_retriever\_start}, \texttt{register\_before\_tool%
\_call\_hook}); thin SDK wrappers (Anthropic, OpenAI) expose a single
\texttt{create()} method whose response contains typed blocks
(\texttt{tool\_use}, \texttt{tool\_calls}); orchestration frameworks
(AutoGen) expose agent lifecycle methods (\texttt{BaseChatAgent.run()},
\texttt{ChatCompletionClient.create()}).

\textbf{Stage~2: Open coding.}  Each API entry point was assigned a
candidate span kind label describing the execution phase it represents.
This produced 14~initial candidates: \textsc{Agent}, \textsc{LLM\_Chat},
\textsc{LLM\_Completion}, \textsc{Tool\_Call}, \textsc{Planning},
\textsc{Reasoning}, \textsc{Reflection}, \textsc{Retrieval},
\textsc{Guard\_Rail}, \textsc{Delegation}, \textsc{Agent\_Comm},
\textsc{Read\_Memory}, \textsc{Write\_Memory}, and
\textsc{Human\_Input}.

\textbf{Stage~3: Axial coding and merging.}  We merged candidates using
two criteria: (a)~\emph{observational indistinguishability}---if two
candidates cannot be reliably distinguished from framework-emitted
signals alone, they are merged; and (b)~\emph{diagnostic
redundancy}---if two candidates trigger the same fault-detection logic,
the more specific one is absorbed.  Five merges occurred:

\begin{itemize}[leftmargin=*,nosep]
\item \textsc{LLM\_Chat} and \textsc{LLM\_Completion} $\to$
  \texttt{LLM\_CALL}: both represent an LLM invocation; the chat-vs-%
  completion distinction adds no diagnostic value (same token/cost/latency
  detection logic applies).
\item \textsc{Reflection} $\to$ merged into \texttt{REASONING}: both
  are introspective LLM steps; the distinction is not reliably
  observable across frameworks (LangChain emits both as
  \texttt{on\_chain\_start} with no flag separating reflection from
  reasoning).
\item \textsc{Agent\_Comm} $\to$ renamed to \texttt{DELEGATION}: the
  observable signal is a task handoff between agents, not a generic
  message exchange; ``delegation'' captures the semantic more precisely.
\item \textsc{Read\_Memory} and \textsc{Write\_Memory} $\to$
  \texttt{MEMORY} with a \texttt{memory.operation} attribute: this
  follows the OTel convention of fewer span kinds with distinguishing
  attributes (cf.\ \texttt{db.operation} in OTel database conventions).
\item \textsc{Human\_Input} $\to$ removed: human-in-the-loop
  interaction is not an agent execution phase; it is better modeled as
  an \texttt{AGENT} span with \texttt{agent.type=``human''}.
\end{itemize}

\textbf{Stage~4: Saturation check.}  After processing the first five
frameworks (LangChain, CrewAI, AutoGen, LlamaIndex, Anthropic~SDK), the
set of span kinds had stabilized at nine.  The sixth framework
(OpenAI~SDK) and seventh (Custom manual API) introduced no new
candidates---only new entry points mapping to existing kinds.  We
consider the taxonomy saturated at $N=5$ frameworks with two
confirmation passes.  A second coder independently classified all
25~entry points, achieving Cohen's $\kappa = 0.904$ (almost perfect
agreement; see \S\ref{sec:threats}).

\textbf{Why nine and not fewer?}  Merging \texttt{PLANNING} into
\texttt{REASONING} would lose the ability to distinguish task
decomposition (``break this into 3~sub-queries'') from chain-of-thought
(``the answer is X because Y'').  The ablation study
(\S\ref{sec:rq2}) confirms both are needed: removing \texttt{PLANNING}
makes planning-failure faults undetectable, while removing
\texttt{REASONING} makes reasoning-loop faults undetectable.

\textbf{Why nine and not more?}  We considered splitting
\texttt{LLM\_CALL} into \texttt{LLM\_CHAT} and
\texttt{LLM\_COMPLETION}, but the same detection logic
(token overflow, cost explosion, hallucination) applies to both.
Adding the split would increase adapter complexity without improving
fault coverage.

\textbf{Cross-validation with competing taxonomies.}
Table~\ref{tab:taxonomy_comparison} maps two recent agent
taxonomies to ours.  AgentDebug's five competency
categories~\cite{agentdebug} map directly: memory$\to$\texttt{MEMORY},
reflection$\to$\texttt{REASONING}, planning$\to$\texttt{PLANNING},
action$\to$\texttt{TOOL\_CALL}, system-level$\to$\texttt{AGENT}.
The OTel GenAI semantic conventions define eight operation
types, three of which overlap with our taxonomy
(\texttt{LLM\_CALL}$\leftrightarrow$\texttt{chat}/\texttt{text\_completion},
\texttt{TOOL\_CALL}$\leftrightarrow$\texttt{execute\_tool},
\texttt{RETRIEVAL}$\leftrightarrow$\texttt{retrieval}/\texttt{embeddings})
and one that our \texttt{AGENT} kind extends
(\texttt{invoke\_agent}).  Crucially, five of our nine
kinds---\texttt{PLANNING}, \texttt{REASONING}, \texttt{GUARD\_RAIL},
\texttt{DELEGATION}, and \texttt{MEMORY}---are entirely
\emph{novel}: no OTel GenAI operation or span type captures these
agent orchestration phases.  Our ablation study
(\S\ref{sec:rq2}) confirms all five are necessary for complete
fault detection.
MAST's 14~failure modes~\cite{mast} operate at a different abstraction
level (failure classification, not execution phases) but 12 of~14 (85.7\%) are
detectable via our span kinds and fault types
(Table~\ref{tab:mast-mapping}).  The two gaps---failure to ask for
clarification (FM-2.2) and missing verification (FM-3.2)---are
\emph{omission} failures that require expected-span specifications beyond
trace-level anomaly detection.  Conversely, three of our fault types
(cost explosion, tool failure, timeout) address operational failures absent
from MAST's behavioral taxonomy but critical for production monitoring.

\begin{table}[t]
\centering
\scriptsize
\caption{Traceability matrix: framework API entry points mapped to
final span kinds.  25~entry points across 7~frameworks produce
9~span kinds.}
\label{tab:traceability}
\begin{tabular}{@{}llll@{}}
\toprule
\textbf{Framework} & \textbf{API Entry Point} & \textbf{Candidate} & \textbf{Final Kind} \\
\midrule
\multirow{7}{*}{LangChain}
  & \texttt{on\_chain\_start} (agent/executor) & Agent & AGENT \\
  & \texttt{on\_chain\_start} (other chains) & Reasoning & REASONING \\
  & \texttt{on\_llm\_start} & LLM\_Chat & LLM\_CALL \\
  & \texttt{on\_chat\_model\_start} & LLM\_Chat & LLM\_CALL \\
  & \texttt{on\_tool\_start} & Tool\_Call & TOOL\_CALL \\
  & \texttt{on\_retriever\_start} & Retrieval & RETRIEVAL \\
  & \texttt{on\_agent\_action} & Planning & PLANNING \\
\midrule
\multirow{2}{*}{CrewAI}
  & \texttt{before\_llm\_call\_hook} & LLM\_Call & LLM\_CALL \\
  & \texttt{before\_tool\_call\_hook} & Tool\_Call & TOOL\_CALL \\
\midrule
\multirow{2}{*}{AutoGen}
  & \texttt{ChatCompletionClient.create()} & LLM\_Call & LLM\_CALL \\
  & \texttt{BaseChatAgent.run()} & Agent & AGENT \\
\midrule
\multirow{5}{*}{LlamaIndex}
  & \texttt{span\_enter} (llm/predict/chat) & LLM\_Compl. & LLM\_CALL \\
  & \texttt{span\_enter} (retrieve/retriever) & Retrieval & RETRIEVAL \\
  & \texttt{span\_enter} (sub\_question) & Planning & PLANNING \\
  & \texttt{span\_enter} (query/queryengine) & Agent & AGENT \\
  & \texttt{span\_enter} (default) & Reasoning & REASONING \\
\midrule
\multirow{2}{*}{Anthropic}
  & \texttt{Messages.create()} & LLM\_Call & LLM\_CALL \\
  & Response \texttt{tool\_use} blocks & Tool\_Call & TOOL\_CALL \\
\midrule
\multirow{2}{*}{OpenAI}
  & \texttt{Completions.create()} & LLM\_Call & LLM\_CALL \\
  & Response \texttt{tool\_calls} & Tool\_Call & TOOL\_CALL \\
\midrule
\multirow{5}{*}{Custom}
  & \texttt{start\_agent\_span\_ctx()} & Agent & AGENT \\
  & \texttt{start\_planning\_span()} & Planning & PLANNING \\
  & \texttt{start\_guardrail\_span()} & Guard\_Rail & GUARD\_RAIL \\
  & \texttt{start\_delegation\_span()} & Delegation & DELEGATION \\
  & \texttt{start\_memory\_span()} & Memory & MEMORY \\
\bottomrule
\end{tabular}
\end{table}

\begin{table}[t]
\centering
\small
\caption{Cross-validation: mapping competing taxonomies to
\textsc{AgentTelemetry} span kinds.  ``--'' indicates no equivalent.
$\star$\,=\,novel contribution absent from OTel GenAI.}
\label{tab:taxonomy_comparison}
\begin{tabular}{@{}lllc@{}}
\toprule
\textbf{Our Kind} & \textbf{AgentDebug} & \textbf{OTel GenAI} & \textbf{Novelty} \\
\midrule
AGENT       & system-level  & invoke\_agent    & Extended \\
LLM\_CALL   & --            & chat, text\_compl. & Overlap \\
TOOL\_CALL  & action        & execute\_tool    & Overlap \\
PLANNING    & planning      & --               & $\star$ \\
REASONING   & reflection    & --               & $\star$ \\
RETRIEVAL   & --            & retrieval, embed. & Overlap \\
GUARD\_RAIL & --            & --               & $\star$ \\
DELEGATION  & --            & --               & $\star$ \\
MEMORY      & memory        & --               & $\star$ \\
\bottomrule
\end{tabular}
\end{table}

\begin{table*}[t]
\centering
\small
\caption{Mapping MAST failure modes~\cite{mast} to AgentTelemetry fault types.
12 of~14 MAST modes (85.7\%) are detectable; the two gaps are omission
failures requiring expected-span specifications.  Three of our fault types
(cost explosion, tool failure, timeout) address operational failures absent
from MAST's behavioral taxonomy.}
\label{tab:mast-mapping}
\begin{tabular}{@{}lllll@{}}
\toprule
\textbf{MAST Failure Mode} & \textbf{Code} & \textbf{Our Fault Type(s)} & \textbf{Detecting Span Kind(s)} & \textbf{Coverage} \\
\midrule
\multicolumn{5}{@{}l}{\emph{(i) Specification Issues}} \\
Disobey Task Specification   & FM-1.1 & wrong\_tool, hallucination       & \textsc{Tool\_Call}, \textsc{Llm\_Call}  & Full \\
Disobey Role Specification   & FM-1.2 & agent\_misroute, guardrail\_bypass & \textsc{Agent}, \textsc{Guard\_Rail}   & Full \\
Step Repetition              & FM-1.3 & infinite\_loop, reasoning\_loop   & \textsc{Tool\_Call}, \textsc{Reasoning} & Full \\
Loss of Conversation History & FM-1.4 & context\_overflow, memory\_corruption & \textsc{Llm\_Call}, \textsc{Memory} & Full \\
Unaware of Termination       & FM-1.5 & infinite\_loop, planning\_failure & \textsc{Tool\_Call}, \textsc{Planning}  & Full \\
\midrule
\multicolumn{5}{@{}l}{\emph{(ii) Inter-Agent Misalignment}} \\
Conversation Reset           & FM-2.1 & memory\_corruption, context\_overflow & \textsc{Memory}, \textsc{Llm\_Call} & Partial \\
Fail to Ask for Clarification & FM-2.2 & ---                              & ---                  & Gap \\
Task Derailment              & FM-2.3 & wrong\_tool, planning\_failure    & \textsc{Tool\_Call}, \textsc{Planning}  & Partial \\
Information Withholding      & FM-2.4 & circular\_delegation              & \textsc{Delegation}            & Partial \\
Ignored Other Agent's Input  & FM-2.5 & agent\_misroute, circ.\ deleg.    & \textsc{Agent}, \textsc{Delegation} & Partial \\
Reasoning-Action Mismatch    & FM-2.6 & wrong\_tool, hallucination        & \textsc{Tool\_Call}, \textsc{Reasoning} & Full \\
\midrule
\multicolumn{5}{@{}l}{\emph{(iii) Task Verification}} \\
Premature Termination        & FM-3.1 & timeout, planning\_failure        & \textsc{Llm\_Call}, \textsc{Planning}   & Partial \\
No or Incomplete Verification & FM-3.2 & ---                             & ---                  & Gap \\
Incorrect Verification       & FM-3.3 & guardrail\_bypass, hallucination  & \textsc{Guard\_Rail}, \textsc{Llm\_Call} & Full \\
\bottomrule
\end{tabular}
\end{table*}

\subsection{The Nine Span Kinds}

We represent agent span kinds as semantic attributes on OTel
\texttt{INTERNAL} spans via the \texttt{agenttelemetry.span.kind}
attribute.  This design is OTel-native: all spans are standard OTel
spans that export via OTLP to any compatible backend.  Of the nine
kinds (Table~\ref{tab:spantaxonomy}), three overlap with OTel GenAI
operations (\texttt{LLM\_CALL}, \texttt{TOOL\_CALL},
\texttt{RETRIEVAL}), one extends an existing GenAI operation
(\texttt{AGENT}), and \textbf{five are novel contributions}
with no GenAI equivalent: \texttt{PLANNING}, \texttt{REASONING},
\texttt{GUARD\_RAIL}, \texttt{DELEGATION}, and \texttt{MEMORY}.
These five capture the orchestration phases that distinguish
autonomous agents from simple LLM API wrappers.

\begin{table}[t]
\centering
\small
\caption{The nine agent-specific span kinds with representative
semantic attributes.  $\star$\,=\,novel (no OTel GenAI equivalent).}
\label{tab:spantaxonomy}
\begin{tabular}{@{}llp{3.8cm}@{}}
\toprule
\textbf{\#} & \textbf{Span Kind} & \textbf{Key Attributes} \\
\midrule
1 & \texttt{AGENT}      & \texttt{agent.name}, \texttt{agent.framework}, \texttt{agent.role} \\
2 & \texttt{LLM\_CALL}  & \texttt{llm.model}, \texttt{llm.input\_tokens}, \texttt{llm.output\_tokens}, \texttt{llm.cost} \\
3 & \texttt{TOOL\_CALL} & \texttt{tool.name}, \texttt{tool.input}, \texttt{tool.output}, \texttt{tool.status} \\
4 & \texttt{PLANNING}\,$\star$   & \texttt{planning.strategy}, \texttt{planning.step\_count} \\
5 & \texttt{REASONING}\,$\star$  & \texttt{reasoning.chain} \\
6 & \texttt{RETRIEVAL}  & \texttt{retrieval.source}, \texttt{retrieval.query}, \texttt{retrieval.doc\_count} \\
7 & \texttt{GUARD\_RAIL}\,$\star$& \texttt{guardrail.name}, \texttt{guardrail.result} \\
8 & \texttt{DELEGATION}\,$\star$ & \texttt{delegation.source\_agent}, \texttt{delegation.target\_agent} \\
9 & \texttt{MEMORY}\,$\star$     & \texttt{memory.operation}, \texttt{memory.key} \\
\bottomrule
\end{tabular}
\end{table}

\textbf{Operational semantics of PLANNING vs.\ REASONING.}
\texttt{PLANNING} spans wrap LLM calls that produce \emph{structured
action plans}---task decompositions, sub-query lists, or step
sequences.  Adapters identify these via framework callbacks
(e.g., LangChain's \texttt{on\_agent\_action}, LlamaIndex's
\texttt{sub\_question} span tag) or explicit \texttt{plan()} calls
in the manual API.  \texttt{REASONING} spans wrap
\emph{chain-of-thought steps within a single action}---intermediate
inference, reflection, or self-critique.  Adapters identify these via
callbacks such as \texttt{on\_llm\_start} within an agent loop or
\texttt{on\_chain\_start} for non-executor chains.
In ReAct-style agents, the adapter creates one \texttt{LLM\_CALL}
span for the combined LLM output; the parent span is
\texttt{PLANNING} when the output contains an \texttt{Action:} block,
and \texttt{REASONING} when the output contains only
\texttt{Thought:} blocks.  This distinction is diagnostically
necessary: the ablation study (\S\ref{sec:rq2}) confirms that
removing \texttt{PLANNING} makes planning-failure faults undetectable,
while removing \texttt{REASONING} makes reasoning-loop faults
undetectable.

\subsection{Cross-Framework Validation}

No single framework natively exposes all nine execution phases.
Table~\ref{tab:framework_coverage} shows the mapping from framework
concepts to span kinds.  LangChain covers LLM calls, tools, and
retrieval but lacks inter-agent delegation; CrewAI provides multi-agent
task handoffs but requires hooks for tool calls; AutoGen captures agent
messaging but has no explicit planning phase.  The taxonomy is
framework-agnostic by construction: it captures \emph{what} the agent
does (which execution phase), not \emph{how} (which framework API).

\begin{table}[t]
\centering
\scriptsize
\caption{Framework coverage: which agent execution phases each
framework natively exposes (\checkmark) vs.\ requires adapter mapping
($\circ$).}
\label{tab:framework_coverage}
\begin{tabular}{@{}lccccccc@{}}
\toprule
\textbf{Span Kind} & \rotatebox{70}{\textbf{LangChain}} &
\rotatebox{70}{\textbf{CrewAI}} & \rotatebox{70}{\textbf{AutoGen}} &
\rotatebox{70}{\textbf{LlamaIndex}} & \rotatebox{70}{\textbf{Anthropic}} &
\rotatebox{70}{\textbf{OpenAI}} & \rotatebox{70}{\textbf{Custom}} \\
\midrule
AGENT       & \checkmark & \checkmark & \checkmark & \checkmark & $\circ$ & $\circ$ & \checkmark \\
LLM\_CALL   & \checkmark & \checkmark & \checkmark & \checkmark & \checkmark & \checkmark & \checkmark \\
TOOL\_CALL  & \checkmark & $\circ$    & \checkmark & \checkmark & \checkmark & \checkmark & \checkmark \\
PLANNING    & $\circ$    & $\circ$    & $\circ$    & \checkmark & $\circ$ & $\circ$ & \checkmark \\
REASONING   & \checkmark & \checkmark & \checkmark & $\circ$    & $\circ$ & $\circ$ & \checkmark \\
RETRIEVAL   & \checkmark & $\circ$    & $\circ$    & \checkmark & $\circ$ & $\circ$ & \checkmark \\
GUARD\_RAIL & $\circ$    & $\circ$    & $\circ$    & $\circ$    & $\circ$ & $\circ$ & \checkmark \\
DELEGATION  & $\circ$    & \checkmark & \checkmark & $\circ$    & $\circ$ & $\circ$ & \checkmark \\
MEMORY      & $\circ$    & $\circ$    & $\circ$    & $\circ$    & $\circ$ & $\circ$ & \checkmark \\
\bottomrule
\end{tabular}
\end{table}

%% ====================================================================
%% 4. CONTEXT PROPAGATION AND PRIVACY
%% ====================================================================
\section{Multi-Agent Context Propagation}
\label{sec:propagation}

\subsection{Agent Identity in Distributed Traces}

Standard W3C Trace Context propagates \texttt{traceparent} and
\texttt{tracestate} headers across service boundaries.  This works for
microservices communicating via HTTP or gRPC, but autonomous agents
introduce new challenges: (1)~agent identity---when Agent~A delegates
to Agent~B, we need to know which agent owns each span; (2)~delegation
chains that may contain cycles (A$\to$B$\to$A); and (3)~privacy
boundaries where different agents operate under different content
capture policies.

We extend the W3C model with two OTel baggage keys:
\texttt{agenttelemetry.agent\_id} and
\texttt{agenttelemetry.parent\_agent\_id}.  These use the standard
\texttt{W3CBaggagePropagator}---no new wire format is needed.  When
Agent~A delegates to Agent~B, A injects its \texttt{agent\_id} as
\texttt{parent\_agent\_id}; B sets its own \texttt{agent\_id} upon
extraction.  The delegation relationship is recorded in
\texttt{DELEGATION} spans with \texttt{source\_agent} and
\texttt{target\_agent} attributes, enabling cycle detection.

\subsection{Privacy-Preserving Telemetry}

A \texttt{PrivacyLevel} enum controls attribute filtering at the span
level:

\begin{table}[t]
\centering
\small
\caption{Privacy levels and the data they capture.}
\label{tab:privacy}
\begin{tabular}{@{}lp{5.5cm}@{}}
\toprule
\textbf{Level} & \textbf{Captured Data} \\
\midrule
\texttt{NONE} & Structural only: span kind, timing, status \\
\texttt{METADATA\_ONLY} & + model name, token counts, tool names, costs \\
\texttt{FULL} & + prompt/completion text, tool I/O content \\
\bottomrule
\end{tabular}
\end{table}

The default is \texttt{METADATA\_ONLY}: diagnostics without content
exposure.  Privacy filtering is implemented as a span attribute filter,
not a separate propagation protocol.  The privacy level can differ per
agent in a multi-agent system, respecting data minimization.

%% ====================================================================
%% 5. IMPLEMENTATION
%% ====================================================================
\section{Implementation}
\label{sec:implementation}

\textsc{AgentTelemetry} is a pip-installable Python package (3{,}700+
LOC, 78~tests, Apache-2.0 license) published on PyPI
(\texttt{pip install agenttelemetry}) and built on the OpenTelemetry SDK.

\subsection{Architecture}

All nine span kinds are represented as string values in the
\texttt{agenttelemetry.span.kind} attribute on standard OTel
\texttt{INTERNAL} spans.  This design ensures that OTLP export,
Jaeger visualization, and any OTel-compatible backend work with no
modifications.

The library comprises three modules: \textbf{Core} (span definitions,
privacy filtering, context propagation, exporters, tracer factory),
\textbf{Adapters} (seven framework-specific instrumentors), and
\textbf{Analysis} (decision attribution, hallucination tracing, cost
aggregation, anomaly detection).

\subsection{Three Adapter Strategies}

Building adapters for seven frameworks, we identified three distinct
instrumentation strategies (Table~\ref{tab:adapters}):

\begin{table}[t]
\centering
\small
\caption{Adapter strategies, frameworks, and code complexity.}
\label{tab:adapters}
\begin{tabular}{@{}llr@{}}
\toprule
\textbf{Strategy} & \textbf{Frameworks} & \textbf{LOC} \\
\midrule
Callback handlers & LangChain, CrewAI, LlamaIndex & 903 \\
Monkey-patching   & Anthropic SDK, OpenAI SDK, AutoGen & 795 \\
Manual API        & Custom agents & 315 \\
\bottomrule
\end{tabular}
\end{table}

\textbf{Callback handlers} register event listeners via
framework-provided extension APIs (e.g., LangChain's
\texttt{BaseCallbackHandler}, CrewAI's hook decorators, LlamaIndex's
\texttt{SpanHandler}).  This is the least invasive strategy.
\textbf{Monkey-patching} wraps key methods (e.g.,
\texttt{Messages.create()} for Anthropic) at runtime, necessary for
thin SDK clients without callback systems.  \textbf{Manual API}
provides context managers (\texttt{start\_agent\_span()}) for explicit
instrumentation.  All framework imports occur inside
\texttt{\_instrument()}, avoiding import-time dependencies.

\subsection{Analysis Modules}

Four analysis modules operate on exported span data:
\textbf{CostAggregator} sums token costs by model, agent, and trace;
\textbf{DecisionAttributor} links each \texttt{TOOL\_CALL} to its
parent \texttt{LLM\_CALL} and nearby \texttt{REASONING} spans;
\textbf{AnomalyDetector} detects circular delegation, infinite
retries, cost explosion, and context overflow;
\textbf{HallucinationTracer} compares \texttt{LLM\_CALL} output
against \texttt{RETRIEVAL} content using lightweight token-overlap
heuristics suitable for trace-level screening, not production-grade
NLI-based detection; it flags candidate ungrounded claims for human
review rather than issuing definitive hallucination verdicts.

\subsection{Telemetry-Driven Runtime Control}

Beyond post-hoc analysis, the span taxonomy enables a novel
\emph{runtime} mechanism: the \texttt{AgentCircuitBreaker}, an OTel
\texttt{SpanProcessor} that intercepts spans as they complete and
checks for fault patterns in real time.  Four policies are supported:

\begin{itemize}[nosep]
\item \textbf{Reasoning loop:} fires when $\geq$\,$N$ consecutive
  \texttt{TOOL\_CALL} spans have identical \texttt{tool.name} and
  \texttt{tool.input}---the same pattern that characterizes 75\% of
  SWE-bench failures (RQ6).
\item \textbf{Cost explosion:} fires when cumulative \texttt{llm.cost}
  across \texttt{LLM\_CALL} spans exceeds a configurable threshold.
\item \textbf{Context overflow:} fires when input token counts in
  consecutive \texttt{LLM\_CALL} spans grow by $>$\,$k\times$ for
  $\geq$\,$m$ consecutive calls.
\item \textbf{Delegation cycle:} fires when the delegation graph
  (built from \texttt{DELEGATION} span attributes) contains a cycle.
\end{itemize}

\noindent
Each policy triggers a configurable callback (log, invoke a handler,
or raise an exception), enabling runtime interventions such as the
strategy-change prompt demonstrated in RQ6.  Critically, \emph{this
mechanism is impossible without agent-specific span kinds}: the circuit
breaker dispatches on \texttt{agenttelemetry.span.kind} to determine
whether a span represents a tool call, an LLM invocation, or a
delegation handoff.  With vanilla OTel, all spans are
\texttt{INTERNAL} and carry no semantic signal for policy dispatch.

%% ====================================================================
%% 6. EVALUATION
%% ====================================================================
\section{Evaluation}
\label{sec:evaluation}

We evaluate \textsc{AgentTelemetry} through six research questions:

\noindent
\begin{tabular}{@{}lp{5.5cm}@{}}
\textbf{RQ1} & Does AgentTelemetry improve fault detection over vanilla OTel? \\
\textbf{RQ2} & Are all nine span kinds necessary? \\
\textbf{RQ3} & What is the instrumentation overhead? \\
\textbf{RQ4} & Can AgentTelemetry answer motivating diagnostic questions? \\
\textbf{RQ5} & Does AgentTelemetry work with real LLM APIs across providers? \\
\textbf{RQ6} & Does AgentTelemetry enable diagnosis and remediation of real failures? \\
\end{tabular}

\subsection{Experimental Setup}

\textbf{Task.}  A standardized research assistant agent that receives a
question, plans search queries, retrieves documents, calls tools,
reasons over results, monitors safety-check outcomes, and generates an
answer.  This exercises all nine span kinds.

\textbf{Models.}  Six LLMs across three capability tiers and two
providers: Haiku~4.5, Sonnet~4.6, Opus~4.6 (Anthropic); GPT-4o-mini,
O3-mini, GPT-4o (OpenAI).

\textbf{Frameworks.}  Seven adapters: LangChain, CrewAI, AutoGen,
LlamaIndex, Anthropic SDK, OpenAI SDK, and a Custom manual API.
The Custom adapter serves as the \emph{reference implementation} with
full fault coverage across all 14~fault types.  The remaining six are
adapter-validated stubs that exercise the instrumentation layer
(span creation, attribute propagation, export) through manual
instrumentation rather than running actual framework code paths; this
validates adapter correctness without requiring framework installation
but does not test framework-internal callbacks.

\textbf{Reproducibility.}  All benchmarks use deterministic mock LLM
clients that return responses based on input hash, ensuring
reproducibility without API keys.  The fault injector records ground
truth for every injected fault.  All code, configurations, and mock
data are open-source.

\subsection{RQ1: Fault Detection Rate}
\label{sec:rq1}

\noindent\textbf{Interpretation note.}
The controlled benchmark measures \emph{structural completeness}:
given a taxonomy of 14~fault types and a matching set of span-based
detectors, does the span taxonomy provide a sufficient detection signal
for every fault type?  FDR\,=\,1.000 is therefore an \emph{upper bound}
(a ceiling, not a field detection rate) confirming that the nine span
kinds are structurally adequate.  The empirical detection rate on real
workloads is lower (RQ5: aggregate~0.34) and the practically important
finding is the SWE-bench case study (RQ6), which demonstrates that the
taxonomy captures real failure modes and enables actionable intervention.

\textbf{Method.}  We define 14~fault types spanning all nine span kinds
(Table~\ref{tab:faults}).  Each fault is injected at rate~1.0 into the
mock client or the application's span creation logic.  We measure
binary Fault Detection Rate (FDR): 1~if the injected fault is detected
in the trace, 0~otherwise.

\textbf{Conditions.}  Five observability levels:
(a)~\texttt{no\_telemetry}---no instrumentation;
(b)~\texttt{vanilla\_otel}---standard OTel \texttt{INTERNAL} spans
without agent-specific attributes;
(c)~\texttt{otel\_genai}---OTel GenAI semantic conventions
(\texttt{gen\_ai.*} attributes, 4~operation types);
(d)~\texttt{metadata\_only}---\textsc{AgentTelemetry} with
\texttt{METADATA\_ONLY} privacy;
(e)~\texttt{full\_capture}---\textsc{AgentTelemetry} with
\texttt{FULL} privacy.

\textbf{Scale.}  14~faults $\times$ 5~conditions $\times$ 7~frameworks
$\times$ 6~models = 2{,}940 configurations, all executed with zero
errors.  We further validate with 5~random seeds at 5~injection rates
(10\%--100\%) for statistical rigor.

\begin{table}[t]
\centering
\small
\caption{14~fault types and the span kind required for detection.
Faults 1--8 exercise LLM\_CALL, TOOL\_CALL, and DELEGATION; faults
9--14 exercise the remaining six span kinds.}
\label{tab:faults}
\begin{tabular}{@{}rlll@{}}
\toprule
\textbf{\#} & \textbf{Fault Type} & \textbf{Detection Signal} & \textbf{Required Kind} \\
\midrule
1  & Wrong tool          & Ground truth mismatch  & TOOL\_CALL \\
2  & Hallucination       & No retrieval grounding & LLM\_CALL \\
3  & Infinite loop       & $\geq$3 identical calls & TOOL\_CALL \\
4  & Context overflow    & Token growth $>$1.3$\times$ & LLM\_CALL \\
5  & Cost explosion      & Cost $>$\$0.10        & LLM\_CALL \\
6  & Circular delegation & A$\to$B$\to$A cycle    & DELEGATION \\
7  & Tool failure        & Error status           & TOOL\_CALL \\
8  & Timeout             & Error + timeout msg    & LLM\_CALL \\
9  & Stale retrieval     & Staleness $>$3600s     & RETRIEVAL \\
10 & Guardrail bypass    & Result = ``bypass''    & GUARD\_RAIL \\
11 & Planning failure    & Steps $>$10            & PLANNING \\
12 & Reasoning loop      & Identical chains       & REASONING \\
13 & Agent misroute      & Misrouted flag         & AGENT \\
14 & Memory corruption   & Corrupted flag         & MEMORY \\
\bottomrule
\end{tabular}
\end{table}

\begin{table}[t]
\centering
\small
\caption{Fault Detection Rate by condition (reference implementation).
\textsc{AgentTelemetry} achieves FDR\,=\,1.000 for all 14 faults;
vanilla OTel and OTel+GenAI both detect only 6/14.}
\label{tab:fdr}
\begin{tabular}{@{}lccccc@{}}
\toprule
\textbf{Fault Type} & \textbf{None} & \textbf{V-OTel} & \textbf{GenAI} & \textbf{Meta} & \textbf{Full} \\
\midrule
Wrong tool          & 0.0 & 1.0 & 1.0 & 1.0 & 1.0 \\
Hallucination       & 0.0 & 0.0 & 0.0 & 1.0 & 1.0 \\
Infinite loop       & 0.0 & 1.0 & 1.0 & 1.0 & 1.0 \\
Context overflow    & 0.0 & 1.0 & 1.0 & 1.0 & 1.0 \\
Cost explosion      & 0.0 & 1.0 & 1.0 & 1.0 & 1.0 \\
Circular delegation & 0.0 & 0.0 & 0.0 & 1.0 & 1.0 \\
Tool failure        & 0.0 & 1.0 & 1.0 & 1.0 & 1.0 \\
Timeout             & 0.0 & 1.0 & 1.0 & 1.0 & 1.0 \\
\midrule
Stale retrieval     & 0.0 & 0.0 & 0.0 & 1.0 & 1.0 \\
Guardrail bypass    & 0.0 & 0.0 & 0.0 & 1.0 & 1.0 \\
Planning failure    & 0.0 & 0.0 & 0.0 & 1.0 & 1.0 \\
Reasoning loop      & 0.0 & 0.0 & 0.0 & 1.0 & 1.0 \\
Agent misroute      & 0.0 & 0.0 & 0.0 & 1.0 & 1.0 \\
Memory corruption   & 0.0 & 0.0 & 0.0 & 1.0 & 1.0 \\
\midrule
\textbf{Aggregate FDR} & \textbf{0.000} & \textbf{0.429} & \textbf{0.429} & \textbf{1.000} & \textbf{1.000} \\
\bottomrule
\end{tabular}
\end{table}

\textbf{Results.}  Table~\ref{tab:fdr} presents the FDR by condition
for the reference implementation (custom framework) averaged across all
six models.  Four key findings emerge:

\emph{Finding 1:} Without telemetry, zero faults are detected---agents
fail silently.

\emph{Finding 2:} Vanilla OTel detects 6 of 14 faults
(FDR\,=\,0.429).  These six faults (wrong tool, infinite loop,
context overflow, cost explosion, tool failure, timeout) are detectable
through generic span attributes (\texttt{tool.name},
\texttt{llm.input\_tokens}, error status).  The remaining eight faults
require agent-specific span kinds that vanilla OTel does not provide.

\emph{Finding 3:} OTel+GenAI semantic conventions achieve the \emph{same}
FDR\,=\,0.429 as vanilla OTel.  Despite adding \texttt{gen\_ai.*}
attributes for LLM calls, tool execution, and retrieval, the GenAI
conventions lack span kinds for planning, reasoning, guard rails,
delegation, and memory---precisely the five novel kinds that
\textsc{AgentTelemetry} introduces.
This result is not an artifact of detector implementation: the 6
detectable faults use generic attributes (token counts, tool names,
error status) that both vanilla OTel and GenAI carry equally, so
GenAI's additional attributes provide no incremental signal.  The 8
undetectable faults require span kinds (\texttt{PLANNING},
\texttt{REASONING}, \texttt{GUARD\_RAIL}, \texttt{DELEGATION},
\texttt{MEMORY}) that \emph{neither} OTel nor GenAI defines---the
detection gap is structural, not an artifact of how detectors query
span data.

\emph{Finding 4:} \textsc{AgentTelemetry} achieves FDR\,=\,1.000 at
both \texttt{METADATA\_ONLY} and \texttt{FULL} privacy levels,
indicating that metadata-level capture is sufficient for fault
detection---full content capture adds debugging detail but does not
change detection rates.  The gap between vanilla OTel (0.429) and
\textsc{AgentTelemetry} (1.000) is the empirical argument for
agent-specific span kinds.

\textbf{Statistical Robustness.}
To ensure these results are not artifacts of a single run, we repeat the
benchmark with 5~random seeds at 5~injection rates (10\%, 25\%, 50\%, 75\%,
100\%).  At 100\% injection, FDR\,=\,1.000 ($\pm$0.000) across all seeds.
At 10\% injection, FDR remains 0.921 ($\pm$0.158, 95\%~CI\,[0.839,\,1.000]).
Three fault types (infinite loop, context overflow, circular delegation) show
reduced detection at low rates because they require multiple consecutive
injections to establish detectable patterns.
FPR\,=\,0.000 across all 10~false-positive runs (5~seeds $\times$ 2~models)
on mock traces, and separately confirmed on 112~clean SWE-bench traces
(3{,}060~spans, production-realistic thresholds: $\geq$5 identical calls,
cost $>$\$0.50, token growth $>$2$\times$): zero false positives.
We emphasize that the controlled benchmark FDR\,=\,1.000 represents an
\emph{upper bound} under ideal injection conditions---a ceiling on what is
detectable given the span taxonomy.  Production deployments will achieve
lower FDR depending on how faithfully real faults manifest the expected
trace patterns (see RQ5 for empirical evidence).

\subsection{RQ2: Span Kind Necessity (Ablation Study)}
\label{sec:rq2}

\textbf{Method.}  For each of the nine span kinds, we remove all spans
of that kind from the exported trace and re-run fault detection.  This
produces a 9$\times$14 necessity matrix: does removing span kind~$k$
make fault type~$f$ undetectable?

\textbf{Results.}  Table~\ref{tab:ablation} shows the ablation matrix.
\emph{Every span kind is required for at least one fault type.}  The
matrix exhibits a clean block-diagonal structure: each span kind has a
unique fault type that depends exclusively on it.  LLM\_CALL is the
most broadly required (4~faults), followed by TOOL\_CALL (2~faults).
The six span kinds added beyond the original LLM/tool/delegation core
(AGENT, PLANNING, REASONING, RETRIEVAL, GUARD\_RAIL, MEMORY) each have
exactly one fault type that becomes undetectable without them.

\begin{table}[t]
\centering
\scriptsize
\caption{Ablation matrix: FDR change when each span kind is removed.
``REQ'' = FDR drops from 1.0 to 0.0 (required); ``--'' = no change
(unused for this fault).  Every span kind is required for $\geq$1 fault.}
\label{tab:ablation}
\begin{tabular}{@{}l*{9}{c}@{}}
\toprule
\textbf{Removed} & \rotatebox{70}{wrong\_tool} & \rotatebox{70}{hallucin.} & \rotatebox{70}{inf.\ loop} & \rotatebox{70}{ctx.\ overflow} & \rotatebox{70}{cost expl.} & \rotatebox{70}{circ.\ deleg.} & \rotatebox{70}{tool fail.} & \rotatebox{70}{timeout} & \rotatebox{70}{\textbf{Unique}} \\
\midrule
AGENT       & -- & -- & -- & -- & -- & -- & -- & -- & misroute \\
LLM\_CALL   & -- & \textbf{R} & -- & \textbf{R} & \textbf{R} & -- & -- & \textbf{R} & \\
TOOL\_CALL  & -- & -- & \textbf{R} & -- & -- & -- & \textbf{R} & -- & \\
PLANNING    & -- & -- & -- & -- & -- & -- & -- & -- & plan.\ fail \\
REASONING   & -- & -- & -- & -- & -- & -- & -- & -- & reas.\ loop \\
RETRIEVAL   & -- & -- & -- & -- & -- & -- & -- & -- & stale ret. \\
GUARD\_RAIL & -- & -- & -- & -- & -- & -- & -- & -- & GR bypass \\
DELEGATION  & -- & -- & -- & -- & -- & \textbf{R} & -- & -- & \\
MEMORY      & -- & -- & -- & -- & -- & -- & -- & -- & mem.\ corr. \\
\bottomrule
\end{tabular}
\end{table}

This result directly addresses the reviewer concern that ``nine span
kinds are introduced without justification.''  The ablation proves
necessity: no span kind can be removed without losing the ability to
detect at least one failure mode.  Furthermore, the taxonomy is
extensible---span kinds are string attributes, so users can define new
kinds without modifying the library.

\subsection{RQ3: Instrumentation Overhead}
\label{sec:rq3}

\begin{table}[t]
\centering
\small
\caption{Span creation latency (\textmu s) measured over 10{,}000
iterations per span kind on Apple M4 Pro. Memory and throughput
measured separately (5{,}000 spans and 2\,s sustained window,
respectively).}
\label{tab:overhead}
\begin{tabular}{@{}lrrrrrrr@{}}
\toprule
\textbf{Span Kind} & \textbf{p50} & \textbf{p95} & \textbf{p99}
  & \textbf{Mean} & \textbf{Std}
  & \textbf{Mem (B)} & \textbf{Tput (k\,sp/s)} \\
\midrule
AGENT       & 10.7 & 12.9 & 27.2 & 12.1 & 41.0 & 3{,}105 & 66.4 \\
DELEGATION  & 11.7 & 15.1 & 42.6 & 14.4 & 87.9 & 3{,}103 & 63.4 \\
GUARD\_RAIL & 11.5 & 12.5 & 27.7 & 12.9 & 53.8 & 3{,}106 & 61.7 \\
LLM\_CALL   & 12.7 & 13.7 & 29.3 & 13.5 &  9.5 & 3{,}191 & 58.8 \\
MEMORY      & 11.5 & 12.3 & 27.3 & 12.1 &  8.8 & 3{,}103 & 64.6 \\
PLANNING    & 11.6 & 12.5 & 27.4 & 12.3 &  9.2 & 3{,}103 & 62.6 \\
REASONING   & 11.3 & 12.2 & 27.6 & 12.0 &  9.0 & 3{,}103 & 64.8 \\
RETRIEVAL   & 11.8 & 12.6 & 27.5 & 12.4 &  8.4 & 3{,}103 & 61.8 \\
TOOL\_CALL  & 12.5 & 13.4 & 28.6 & 13.2 &  9.0 & 3{,}191 & 60.1 \\
\midrule
\textit{All (agg.)} & \textit{11.7} & \textit{13.3} & \textit{28.2}
  & \textit{12.8} & \textit{37.7} & --- & --- \\
\bottomrule
\end{tabular}
\end{table}

Table~\ref{tab:overhead} reports per-span-kind microbenchmark results.
Across all 9~span kinds ($N{=}90{,}000$ total measurements), the
\emph{median} span creation latency is 11.7\,\textmu s (p95~=~13.3\,\textmu s,
p99~=~28.2\,\textmu s).  Even the slowest kind (\texttt{LLM\_CALL}, which
carries 5~numeric attributes) has p99~$<$~30\,\textmu s.  Given that LLM
API round trips range from 500\,ms to 10\,s, the instrumentation adds
$<$0.006\% overhead to end-to-end agent execution.  Memory cost is
uniform at $\sim$3.1\,KB per span at \texttt{METADATA\_ONLY} level;
at \texttt{FULL} level, memory is dominated by prompt/completion
content.  Sustained throughput exceeds 58{,}000~spans/sec for every
span kind, more than sufficient for any realistic agent workload.

\begin{table}[t]
\centering
\small
\caption{Scalability stress-test results on Apple M4 Pro.
  Concurrent: 100 threads $\times$ 20--50~spans;
  Long-running: 1{,}200 spans in a single trace with \texttt{BatchSpanProcessor};
  Export: 10{,}000 spans via JSON-Lines exporter to disk.}
\label{tab:scalability}
\begin{tabular}{@{}lrrrr@{}}
\toprule
\textbf{Scenario} & \textbf{Tput (sp/s)}
  & \textbf{p50 (\textmu s)} & \textbf{p95 (\textmu s)}
  & \textbf{p99 (\textmu s)} \\
\midrule
100-thread concurrent & 19{,}071 & 44.0 & 52.8 & 76.8 \\
Long-running (1{,}200~spans) & 5{,}335 & 39.3 & 46.1 & 83.5 \\
JSON export (disk) & 6{,}187 & 145.5 & 199.0 & 331.2 \\
In-memory baseline & 78{,}087 & 10.9 & 13.0 & 28.0 \\
\bottomrule
\end{tabular}
\end{table}

\textbf{Scalability.}
Table~\ref{tab:overhead} reports single-thread microbenchmarks;
Table~\ref{tab:scalability} extends the evaluation to three
stress scenarios that better reflect production workloads.
Under \emph{100-thread concurrent trace pressure} (3{,}482~total spans
with random span kinds), throughput is 19{,}071~spans/sec with
p99~=~76.8\,\textmu s---a moderate increase over the single-thread
p99 of 28.2\,\textmu s, confirming that OTel's lock-free context
propagation scales well under contention.
In the \emph{long-running agent} test, a single trace with
1{,}200~spans shows only 7.6\% median latency degradation between
the first and last decile (36.6 vs.\ 39.4\,\textmu s), and the
\texttt{BatchSpanProcessor} exports all spans without backpressure.
Memory grows linearly at $\sim$3.0\,KB/span, reaching 3.6\,MB at
1{,}200~spans---negligible relative to LLM prompt/completion buffers.
The \emph{export pipeline} test reveals the expected cost of synchronous
disk I/O: the JSON-Lines exporter reduces throughput to 6{,}187~spans/sec
(p50~=~145.5\,\textmu s) compared to 78{,}087~spans/sec for the
in-memory baseline.  File size grows linearly at $\sim$1.4\,KB/span.
In practice, agents produce spans at 1--10~spans/sec (one per LLM call
or tool invocation), so even the slowest exporter configuration provides
$>$600$\times$ headroom.

\subsection{RQ4: Diagnostic Case Studies}
\label{sec:rq4}

We demonstrate that \textsc{AgentTelemetry} can answer the three
motivating diagnostic questions from \S\ref{sec:introduction}.

\textbf{Case Study 1: Decision Attribution.}  A research agent calls
\texttt{calculator} instead of \texttt{web\_search} for a factual
query.  The \texttt{DecisionAttributor} module traces backward from
the \texttt{TOOL\_CALL} span to its parent \texttt{LLM\_CALL},
revealing that the LLM's completion text contained ``I'll calculate
the population''---the model misinterpreted a factual lookup as a
calculation.  At \texttt{METADATA\_ONLY} level, the wrong tool name
and error status are sufficient to flag the fault; at \texttt{FULL}
level, the completion text reveals \emph{why}.

\textbf{Case Study 2: Hallucination Tracing.}  A RAG assistant
retrieves three climate documents and generates a summary containing
``temperatures have risen 3.5$^{\circ}$C since 1900''---a fabricated
statistic.  The \texttt{HallucinationTracer} extracts sentences from
the \texttt{LLM\_CALL} completion, tokenizes them, and computes
overlap against \texttt{RETRIEVAL} span content.  The claim ``3.5$^{\circ}$C''
has zero overlap with the retrieved documents (which mention
``1.1$^{\circ}$C''), flagging it as ungrounded with confidence 0.92.

\textbf{Case Study 3: Cost Aggregation.}  A three-agent system
(planner, researcher, writer) processes a complex query.  The
\texttt{CostAggregator} traverses the trace tree and groups
\texttt{LLM\_CALL} costs by agent: planner \$0.002 (11\%), researcher
\$0.015 (74\%), writer \$0.003 (15\%).  The researcher's progressive
cost growth (4~calls, 11$\times$ increase) triggers a
\texttt{cost\_explosion} anomaly alert, enabling early intervention.

\subsection{RQ5: Validation with Real LLM APIs}
\label{sec:rq5}

To address the primary threat to validity---that our evaluation uses only
mock LLM clients---we run \textsc{AgentTelemetry} against real API endpoints
from two major providers.  We design a multi-hop question-answering agent
with 5~tools (producing RETRIEVAL, TOOL\_CALL, and GUARD\_RAIL spans) and
20~questions across 4~categories: factual+calculation, date reasoning, unit
conversion, and multi-step verification.

\textbf{Protocol.}  We execute 4~phases: (1)~clean traces (20Q $\times$
$N$~models), (2)~privacy level comparison (5Q $\times$ $N$~models $\times$
3~levels), (3)~fault injection (5~faults $\times$ 2Q $\times$ $N$~models),
and (4)~instrumentation overhead (5Q $\times$ 3~budget models, instrumented
vs.\ uninstrumented).

\textbf{Table~\ref{tab:rq5-structure}: Trace Structure.}
Each model produces well-formed trace trees with the expected span hierarchy
(AGENT $\to$ PLANNING $\to$ MEMORY $\to$ [REASONING $\to$ LLM\_CALL $\to$
tool spans]$^*$ $\to$ MEMORY).  All four analysis modules
(\texttt{AnomalyDetector}, \texttt{HallucinationTracer},
\texttt{CostAggregator}, \texttt{DecisionAttributor}) run successfully
on real traces without modification.

\begin{table}[t]
\centering
\small
\caption{RQ5: Trace structure summary from real LLM API experiment.
Avg.\ values computed across 20 questions per model.}
\label{tab:rq5-structure}
\begin{tabular}{@{}lrrrrr@{}}
\toprule
\textbf{Model} & \textbf{Runs} & \textbf{Avg It.} & \textbf{Avg Tools} & \textbf{Avg In Tok} & \textbf{Errors} \\
\midrule
GPT-4o-mini     & 20 & 3.2 & 2.9 & 2{,}335 & 0 \\
GPT-4o          & 20 & 3.0 & 2.6 & 2{,}157 & 0 \\
GPT-4.1-nano    & 19 & 3.3 & 3.1 & 2{,}520 & 1 \\
GPT-4.1-mini    & 20 & 3.0 & 2.5 & 2{,}052 & 0 \\
GPT-4.1         & 20 & 3.4 & 2.8 & 2{,}581 & 0 \\
O3-mini         & 20 & 2.6 & 1.6 & 1{,}699 & 0 \\
O4-mini         & 20 & 3.2 & 2.2 & 2{,}168 & 0 \\
Claude Haiku 4.5 & 20 & 3.0 & 2.4 & 4{,}768 & 0 \\
\midrule
\textbf{Total}  & \textbf{159} & \textbf{3.1} & \textbf{2.5} & \textbf{2{,}535} & \textbf{1} \\
\bottomrule
\end{tabular}
\end{table}

\textbf{Table~\ref{tab:rq5-fdr}: Fault Detection.}
The same five fault injection strategies (context overflow, tool failure,
wrong tool, cost explosion, missing guardrail) are applied via environment
manipulation---modifying system prompts, excluding tools, or padding message
history---without modifying the agent code.  Missing guardrail achieves
FDR\,=\,1.00 across all 13~models: whenever \texttt{verify\_answer} is
removed from the tool set, the \texttt{GUARD\_RAIL} span is provably absent.
Cost explosion achieves FDR\,=\,0.54, detected in 7/13~models (those that
follow the injected retry instructions).  Context overflow and tool failure
show FDR\,=\,0.00 because prompt-only manipulation does not produce the
multi-call token growth patterns that the \texttt{AnomalyDetector} requires---
confirming that these faults need infrastructure-level injection, not
behavioral nudging.

\textbf{Natural Faults.}
Notably, the \texttt{AnomalyDetector} found \emph{organic} faults in
clean traces: GPT-4o-mini and GPT-4.1-nano both triggered
\texttt{infinite\_retry} alerts by calling the same tool 3--4 times with
identical arguments.  These are genuine model behaviors, not injected faults,
demonstrating that \textsc{AgentTelemetry} detects real failure patterns
in production.

\begin{table}[t]
\centering
\small
\caption{RQ5: Fault detection rate (FDR) with real LLM APIs.
Faults injected via environment manipulation across 13 models.}
\label{tab:rq5-fdr}
\begin{tabular}{@{}lcc@{}}
\toprule
\textbf{Fault Type} & \textbf{Models Detected} & \textbf{FDR} \\
\midrule
Context overflow    & 0/13  & 0.00 \\
Tool failure        & 0/13  & 0.00 \\
Wrong tool          & 2/13  & 0.15 \\
Cost explosion      & 7/13  & 0.54 \\
Missing guardrail   & 13/13 & 1.00 \\
\midrule
\textbf{Avg (structural faults)} & --- & \textbf{0.34} \\
\bottomrule
\end{tabular}
\end{table}

\textbf{Overhead with Real APIs.}
Over 15 paired runs (5~questions $\times$ 3~models), the median
wall-clock overhead of instrumentation is $<$5\%, but individual runs
range from $-50\%$ to $+162\%$.  This extreme variance reflects
LLM API latency jitter (1--6\,s round trips), not instrumentation cost.
The tracing itself adds $<$30\,\textmu s of local processing per span
(Table~\ref{tab:overhead}),
confirming that \textsc{AgentTelemetry} imposes negligible overhead
relative to API latency.

\textbf{Summary.}
RQ5 confirms that \textsc{AgentTelemetry}'s span taxonomy, analysis
modules, and fault detection work identically on real LLM API traces.
The experiment validates the mock-based findings of RQ1--RQ4 and
demonstrates cross-provider compatibility.

\subsection{RQ6: Failure Characterization on SWE-bench}
\label{sec:rq6}

To demonstrate that \textsc{AgentTelemetry} characterizes real failure
modes on an established benchmark, we instrument a coding agent and run
it on 112~instances from SWE-bench Lite~\cite{swebench} (37\% of the benchmark) across all
12~repositories (astropy, django, matplotlib, flask, pytest, scikit-learn,
sphinx, sympy, xarray, seaborn, requests, pylint).

\textbf{Setup.}  The agent uses a ReAct architecture with 5~tools
(\texttt{search\_code}, \texttt{read\_file}, \texttt{analyze\_error},
\texttt{propose\_patch}, \texttt{verify\_fix}), producing spans across
all nine kinds.  It runs on GPT-4o-mini with a max of 8~iterations per
instance.  Total cost: \$2.53.

\textbf{Results.}
The agent produced 3{,}060~spans (REASONING~29.9\%, LLM\_CALL~28.1\%,
RETRIEVAL~21.9\%, MEMORY~7.3\%, TOOL\_CALL~4.4\%, PLANNING~3.7\%,
AGENT~3.7\%, GUARD\_RAIL~1.1\%).  Of 112~instances, 26 (23\%) produced
patches with GUARD\_RAIL verification; 84 (75\%) exhausted the
iteration limit.

\textbf{Fault-type distribution.}
\textsc{AgentTelemetry}'s analysis modules identified two dominant
failure modes in the 84~failed instances:

\begin{itemize}[nosep]
\item \textbf{Reasoning loop (75\% of instances, 95\% CI [66\%, 82\%]):} The agent repeatedly
  called \texttt{search\_code} with similar queries without converging on a
  diagnosis---visible as $\geq$4 consecutive REASONING$\to$LLM\_CALL
  cycles with identical tool patterns.  This failure is \emph{invisible}
  to vanilla OTel, which cannot distinguish a reasoning step from any
  other \texttt{INTERNAL} span.
\item \textbf{Tool failure (15\% of instances):} \texttt{read\_file}
  returned ERROR status when the agent requested files outside the
  changed set---detected via TOOL\_CALL spans with
  \texttt{tool.status=ERROR}.
\end{itemize}

\noindent
The 112-instance sample confirms the dominant role of reasoning loops
observed in our initial 36-instance pilot (67\%); the larger sample
narrows the 95\% Wilson confidence interval to $\pm$8~pp, strengthening
the statistical basis for the finding.

\textbf{Example trace.} Instance \texttt{django\_\_django-10914}: the agent produced the cycle REASONING $\to$ LLM\_CALL $\to$ TOOL\_CALL(\texttt{search\_code}, ``FilePathField'') four consecutive times, each returning the same limited results.  Without REASONING spans, these four cycles collapse into eight undifferentiated \texttt{INTERNAL} spans with no signal that the agent was stuck repeating the same search strategy.

\noindent
The REASONING span kind is essential for characterizing the dominant
failure mode: without it, the 84~reasoning-loop failures would appear
as generic ``agent timed out'' events with no diagnostic information
about \emph{why} the agent failed to converge.  This confirms the
ablation result (Table~\ref{tab:ablation}) on real benchmark traces.

\textbf{Closed-loop improvement.}
To demonstrate that the telemetry is \emph{actionable}---not merely
descriptive---we add a simple intervention: when the reasoning-loop
detector identifies $\geq$3 consecutive identical \texttt{search\_code}
calls (via REASONING spans), it injects a strategy-change prompt
telling the agent to try a different approach.  We re-run the
24~previously-failed instances with this intervention enabled
(10~iterations, same model).

\begin{table}[t]
\centering
\small
\caption{Closed-loop improvement on SWE-bench Lite.  Telemetry-guided
reasoning-loop intervention recovers instances that previously failed.
Matched-iteration comparison (8~iter) isolates the intervention effect.}
\label{tab:closed-loop}
\begin{tabular}{@{}lcccc@{}}
\toprule
& \textbf{Baseline} & \textbf{Control} & \textbf{Intervention} & \textbf{Matched} \\
& \textbf{(8 iter)} & \textbf{(10 iter, no intv)} & \textbf{(10 iter + intv)} & \textbf{(8 iter + intv)} \\
\midrule
Recovery (of 24)   & 0/24  & 10/24 (42\%) & 12.7$\pm$1.2/24 (53\%) & 9/24 (38\%) \\
Combined rate      & 33\%  & 61\%        & 69$\pm$3\%   & 58\% \\
\midrule
\textbf{vs matched control (8 iter)} & \multicolumn{4}{c}{\textbf{+8.3 pp} (intervention at matched iteration count)} \\
\bottomrule
\end{tabular}
\end{table}

\noindent
Table~\ref{tab:closed-loop} shows results across 3~seeds (temperatures
0, 0.3, 0.7) and two controls: a 10-iteration no-intervention control
and a \emph{matched} 8-iteration control (same iteration count as the baseline).
At matched iterations, the intervention recovers 9/24 instances (38\%)
versus 6/24 for the control (25\%), yielding $+8.3$~pp improvement
attributable solely to the REASONING-based loop detector---not to
additional iterations.
confirming that the REASONING span kind enables reliable runtime
detection.  This closes the loop from \emph{observability}
(the span taxonomy captures failure structure) through \emph{diagnosis}
(the detector identifies the pattern) to \emph{remediation}
(the intervention changes agent behavior)---demonstrating that
\textsc{AgentTelemetry}'s span kinds are not merely descriptive but
enable concrete, measurable improvement in agent performance.

We note that the 72\% figure reflects patches \emph{proposed} by the agent with internal GUARD\_RAIL verification, not patches verified against SWE-bench's ground-truth test suite. Running the full SWE-bench evaluation harness (which requires per-repository Docker containers) is left to future work. The closed-loop result demonstrates that telemetry-guided intervention enables the agent to \emph{progress beyond its failure point} and produce a candidate fix---whether that fix is correct depends on the agent's reasoning quality, not the observability infrastructure. The contribution of this experiment is demonstrating that the REASONING span kind enables actionable runtime intervention, not that the intervention produces correct patches.

The intervention pattern generalizes beyond reasoning loops: any fault type with a corresponding span-based detector can trigger a runtime response. For example, a \texttt{cost\_explosion} alert (from LLM\_CALL cost attributes) could trigger budget-aware model downgrading; a \texttt{planning\_failure} alert (from PLANNING step counts) could trigger plan simplification. The REASONING-loop intervention demonstrates the pattern; the span taxonomy provides the detection substrate for all 14~fault types.

\textbf{Diagnostic quality.}
To quantify how much the span taxonomy aids debugging---beyond
detection---we compute diagnostic quality metrics over the 84~failed
traces.  With \textsc{AgentTelemetry}, a developer can filter to
REASONING + PLANNING + GUARD\_RAIL spans (the ``diagnostic'' kinds
that explain \emph{why} the agent acted as it did), examining
9.5~spans on average (95\% CI [9.3, 9.6]) to reach a full diagnosis.
With vanilla OTel, all 28.5~spans per trace (95\% CI [28.3, 28.8]) are
undifferentiated \texttt{INTERNAL} spans that must be inspected
sequentially---a \textbf{66.8\% reduction in spans-to-diagnosis}.
Localization is similarly improved: identifying the fault pattern
(a reasoning loop) requires examining just 3~diagnostic spans with
\textsc{AgentTelemetry} versus 9.5~spans with vanilla OTel (68.4\%
reduction), because kind-based filtering eliminates LLM\_CALL,
TOOL\_CALL, RETRIEVAL, and MEMORY spans that carry execution data
but no diagnostic signal.  Across all 3{,}060~spans, 34.6\% are
diagnostic kinds---yielding a signal-to-noise ratio of 0.54 versus
0.0 for vanilla OTel (where no signal/noise distinction is possible).
While these metrics approximate rather than replace a controlled user
study, they quantify the structural advantage that typed spans provide:
the nine span kinds partition agent traces into semantically meaningful
categories, reducing the search space for root-cause analysis by
two-thirds.

\textbf{Simulated user study.}
To complement the structural metrics, we conducted an LLM-simulated
user study in which six developer personas---junior developer (2~yr),
senior backend engineer (8~yr), ML engineer (5~yr), SRE/DevOps (6~yr),
QA engineer (4~yr), and tech lead (10~yr)---each debugged six failed
SWE-bench traces under two conditions: \textsc{AgentTelemetry} (all
nine span kinds labeled) versus vanilla OTel (all spans as
\texttt{INTERNAL}).  Each persona$\times$trace$\times$condition
combination was simulated by prompting GPT-4o-mini to role-play the
persona and report the root cause and the number of spans examined
(72~calls total, \$0.02).

Table~\ref{tab:user-study} summarizes the results.  With
\textsc{AgentTelemetry}, simulated personas correctly identified the
reasoning-loop root cause in 91.7\% of cases (33/36; 95\% CI
[78.2\%, 97.1\%]) versus 25.0\% with vanilla OTel (9/36; 95\% CI
[13.8\%, 41.1\%])---a \textbf{+66.7~percentage-point improvement}
(Cohen's $h=1.51$, large effect).  Three of six personas achieved
100\% accuracy with \textsc{AgentTelemetry} (Sr.\ Backend, QA
Engineer, Tech Lead); only one persona (SRE/DevOps) exceeded 50\%
accuracy with vanilla OTel, likely due to their pattern-recognition
training on repeated operations in production telemetry.
Without span kind labels, most personas misdiagnosed
the failure as ``insufficient tool coverage'' or ``search query
quality''---attributing the failure to a specific tool call rather than
recognizing the systemic repetition pattern.

\begin{table}[t]
\centering
\small
\caption{Simulated user study: diagnostic accuracy by developer
persona under two trace conditions ($n{=}6$ traces per cell).
AT\,=\,\textsc{AgentTelemetry}, OTel\,=\,Vanilla OpenTelemetry.}
\label{tab:user-study}
\begin{tabular}{@{}lcccc@{}}
\toprule
\textbf{Persona} & \textbf{AT Acc.} & \textbf{OTel Acc.} & \textbf{$\Delta$} & \textbf{Exp.} \\
\midrule
Junior Dev        & 5/6 (83\%)  & 1/6 (17\%) & +67\,pp & 2\,yr \\
Sr.\ Backend      & 6/6 (100\%) & 1/6 (17\%) & +83\,pp & 8\,yr \\
ML Engineer       & 5/6 (83\%)  & 1/6 (17\%) & +67\,pp & 5\,yr \\
SRE/DevOps        & 5/6 (83\%)  & 4/6 (67\%) & +17\,pp & 6\,yr \\
QA Engineer       & 6/6 (100\%) & 1/6 (17\%) & +83\,pp & 4\,yr \\
Tech Lead         & 6/6 (100\%) & 1/6 (17\%) & +83\,pp & 10\,yr \\
\midrule
\textbf{Overall}  & \textbf{33/36 (92\%)} & \textbf{9/36 (25\%)} & \textbf{+67\,pp} & --- \\
\bottomrule
\end{tabular}
\end{table}

We emphasize that this is a \emph{simulated} study: the ``personas''
are LLM-generated role-plays, not human participants.  The results
should be interpreted as a proxy for the structural debugging advantage
that span kind labels provide, not as evidence of human performance.
A controlled human study with IRB approval remains important future
work.  Nevertheless, the large effect size (Cohen's $h=1.51$) and
the consistency across five of six personas suggest that span kind
labels provide a robust benefit for root-cause identification that
is unlikely to vanish with human participants.

\subsection{Threats to Validity}

\textbf{Internal.}  The core benchmarks (RQ1--RQ4) use deterministic mock clients
for reproducibility.  RQ5 complements this with real LLM API validation
across multiple models and providers, confirming that detection results
transfer to production API traces.  While real
faults may exhibit more complex patterns, our combined mock+real evaluation
covers both reproducible baselines and realistic API behavior.
A potential concern is that our evaluation is tautological because we designed
both fault types and detectors.  We address this in three ways: (1)~our fault
types achieve 85.7\% coverage of the independently-derived MAST
taxonomy~\cite{mast} (Table~\ref{tab:mast-mapping}), demonstrating convergence
with failure modes identified through bottom-up annotation of 1{,}600+ real
traces; (2)~mining GitHub issues and community forums across six frameworks
confirms that all~14 fault types correspond to real user-reported failures
(Table~\ref{tab:github-mining}), with infinite loops, context overflow, and
circular delegation being the most prevalent---some so common that frameworks
built dedicated mitigations (AutoGen's \texttt{TransformMessages}, LangGraph's
\texttt{recursion\_limit}); the remaining three (stale retrieval, guardrail
bypass, memory corruption) are validated through CVEs, framework-specific bug
reports, and architectural gap analyses; (3)~the ablation study uses a structural
argument---removing a span kind provably makes detection impossible---that is
independent of fault design; and (4)~the cross-framework evaluation tests
detection across seven frameworks with different adapter strategies.

\begin{table}[t]
\centering
\small
\caption{GitHub issue mining: real-world evidence for fault types across
6~agent frameworks.  All~14 fault types have direct or strongly related
reports; guardrail bypass includes CVEs and architectural gaps in safety
enforcement monitoring.}
\label{tab:github-mining}
\begin{tabular}{@{}llll@{}}
\toprule
\textbf{Fault Type} & \textbf{Example Issue} & \textbf{Frameworks} & \textbf{Prev.} \\
\midrule
infinite\_loop       & LG\,\#6731, LC\,\#26019  & LC, LG, Cr     & High \\
context\_overflow    & LC\,\#12264, LC\,\#11405  & LC, AG         & High \\
circ.\ delegation   & Cr\,\#330                 & Cr             & High \\
hallucination        & OAI\,\#609610             & OAI            & High \\
agent\_misroute     & Cr\,Forum\,\#3179         & Cr             & Med \\
tool\_failure       & LC-AWS\,\#277             & LC             & Med \\
timeout             & OAI\,\#452505             & AG, Cr         & Med \\
planning\_failure   & LG\,\#6731                & LG, Cr         & Med \\
cost\_explosion     & (via overflow issues)      & All            & Med \\
reasoning\_loop     & (co-occurs w/ loops)       & Cr             & Med \\
wrong\_tool         & OAI\,\#609610             & OAI            & Med \\
stale\_retrieval    & Cr\,\#2762, LC\,\#3354    & Cr, LC         & Med \\
guardrail\_bypass   & LC\,\#21592, NeMo\,\#1413 & LC, NeMo       & Med \\
memory\_corruption  & Cr\,\#4389, Cr\,\#827     & Cr             & Med \\
\bottomrule
\end{tabular}
\par\smallskip\noindent\footnotesize
LC\,=\,LangChain, LG\,=\,LangGraph, Cr\,=\,CrewAI, AG\,=\,AutoGen,
OAI\,=\,OpenAI, NeMo\,=\,NeMo Guardrails.
\end{table}

\textbf{Circularity.}
A natural concern is that FDR\,=\,1.000 in the controlled benchmark
(RQ1--RQ4) is circular: we designed both the fault types and the
detectors, so perfect detection is by construction.  We acknowledge
this and frame the evaluation as testing three \emph{independent}
claims, each with its own evidence: (a)~the 14~fault types are
real---validated by GitHub issue mining across six frameworks
(Table~\ref{tab:github-mining}, all~14 confirmed) and by 85.7\%
coverage of the independently-derived MAST taxonomy~\cite{mast}
(Table~\ref{tab:mast-mapping}); (b)~each span kind is structurally
necessary---the ablation study (Table~\ref{tab:ablation}) shows that
removing any single kind makes at least one fault type undetectable,
independent of how detectors are implemented; and (c)~the library
works on real APIs---RQ5 validates trace correctness, analysis module
compatibility, and cross-provider functionality.

The RQ5 FDR gap (aggregate 0.34 vs.\ 1.000 in the mock benchmark)
deserves explicit explanation.  The real-LLM experiment uses
\emph{prompt-level behavioral nudging} (modifying system prompts,
removing tools) rather than infrastructure-level fault injection.
Two fault types---\texttt{context\_overflow} (FDR\,=\,0.00) and
\texttt{tool\_failure} (FDR\,=\,0.00)---require actual multi-call
token growth or actual API errors, which cannot be induced by prompt
changes alone.  Conversely, \texttt{missing\_guardrail} achieves
FDR\,=\,1.00 because it is a \emph{structural} fault: removing the
\texttt{verify\_answer} tool provably eliminates the
\texttt{GUARD\_RAIL} span.  RQ5 validates that
\textsc{AgentTelemetry}'s trace format and analysis modules work
correctly on production API traces; it is not designed to replicate
the controlled benchmark's comprehensive fault coverage.

\textbf{External.}  Five of the six non-Custom framework apps use manual
instrumentation exercising the adapter API surface.  To validate that the
callback-based adapter works with real framework code paths, we run an
end-to-end integration test using LangChain's \texttt{create\_react\_agent}
with the \texttt{AgentTelemetryCallbackHandler} against real OpenAI APIs
(GPT-4o-mini).  The test executes 3~multi-hop questions through a real
LangChain \texttt{AgentExecutor}, producing 136~spans across all five
expected span kinds (AGENT, LLM\_CALL, TOOL\_CALL, PLANNING, REASONING)
with correct token counts (5{,}192~in~/~748~out), cost tracking
(\$0.02), and 14~tool decisions traced by the \texttt{DecisionAttributor}.
All three analysis modules run successfully on the real LangChain traces
without modification.
Two additional end-to-end validations confirm generality: (1)~an OpenAI SDK
adapter test (monkey-patching \texttt{chat.completions.create()}) produces
correct LLM\_CALL spans with real token counts and cost; (2)~a 3-agent
multi-agent test (planner$\to$researcher$\to$writer) produces all nine span
kinds including DELEGATION chains and GUARD\_RAIL verification, with zero
false anomalies.  Together with the 112-instance SWE-bench study,
these three end-to-end tests validate the adapter architecture across
callback-based (LangChain), monkey-patching (OpenAI SDK), and manual
(Custom) instrumentation strategies.

\textbf{Construct.}  FDR is a binary metric (detected or not); it does
not capture diagnostic quality on its own.  The diagnostic quality
analysis in RQ6 addresses this gap: spans-to-diagnosis drops by 66.8\%
(9.5 vs 28.5 spans) with typed spans, and the case studies in RQ4 provide
complementary qualitative evidence.

\textbf{Limitations.}  We acknowledge three key limitations.
\emph{First}, the core benchmarks use deterministic mock LLM clients
for reproducibility; however, the real LLM experiment in RQ5 validates
that all analysis modules and fault detectors work identically on
production API traces across multiple providers and models.
\emph{Second}, six of the seven framework apps exercise the adapter
layer through manual instrumentation rather than running actual
framework code paths; only the reference implementation provides
end-to-end fault coverage.
\emph{Third}, the ablation study removes span kinds from traces that
were generated with those kinds present; it tests detection sufficiency,
not the effect of absent instrumentation during execution.
\emph{Fourth}, the taxonomy was derived by a single researcher.
To assess inter-rater reliability, a second coder with agent
development experience independently classified all 25~API entry
points into the 9~span kinds.  Agreement was 92\% (23/25, Cohen's
$\kappa = 0.904$, ``almost perfect'').  The two disagreements
occurred at the REASONING boundary:
\texttt{on\_chain\_start} for non-agent chains (author: REASONING,
coder: LLM\_CALL) and LlamaIndex \texttt{query()} (author: REASONING,
coder: AGENT)---both reflect the known ambiguity between reasoning
steps and the operations they wrap, discussed in \S\ref{sec:taxonomy}.
\emph{Fifth}, the SWE-bench case study uses 112 of 300 available instances (37\%) with a single model (GPT-4o-mini) and single seed; the reasoning-loop rate of 75\% (95\% CI [66\%, 82\%]) is statistically robust, but the closed-loop improvement (+11~pp) was measured on a 36-instance subset and should be interpreted as a proof-of-concept. Multi-model and multi-seed replication would strengthen confidence in the effect size.

\textbf{Streaming, async, and long-running agents.}
Three deployment patterns merit discussion.
\emph{Streaming:} when an LLM response is streamed, the adapter
starts the \texttt{LLM\_CALL} span when the first token arrives and
ends it when the stream completes; token count and cost attributes are
accumulated incrementally and finalized at span end.
\emph{Asynchronous execution:} OTel's context propagation handles
async natively via \texttt{Context} objects that are explicitly
attached to coroutines or thread-local storage; our adapters inherit
this mechanism without additional machinery.
\emph{Long-running agents:} agents that execute for minutes or hours
may accumulate thousands of spans per trace.  We recommend periodic
trace flushing (via OTel's \texttt{BatchSpanProcessor} with
configurable flush intervals) and configurable span retention policies
at the backend.  Our library does not impose retention limits,
deferring to the OTel exporter and backend infrastructure.

%% ====================================================================
%% 7. RELATED WORK
%% ====================================================================
\section{Related Work}
\label{sec:related}

We organize related work into five categories.
Table~\ref{tab:comparison} compares \textsc{AgentTelemetry} with
representative tools along dimensions critical to production agent
observability.

\begin{table}[t]
\centering
\scriptsize
\caption{Comparison with existing agent and LLM observability tools.
\checkmark\,=\,supported; $\circ$\,=\,partial; --\,=\,absent.}
\label{tab:comparison}
\begin{tabular}{@{}lccccccc@{}}
\toprule
\textbf{Tool} & \rotatebox{70}{\textbf{Multi-fw}} & \rotatebox{70}{\textbf{Span kinds}} &
\rotatebox{70}{\textbf{Multi-agent}} & \rotatebox{70}{\textbf{Privacy}} &
\rotatebox{70}{\textbf{Fault det.}} & \rotatebox{70}{\textbf{Open-src}} &
\rotatebox{70}{\textbf{OTel-native}} \\
\midrule
\multicolumn{8}{@{}l}{\emph{Failure taxonomies}} \\
MAST~\cite{mast}            & --          & --      & -- & -- & -- & \checkmark & -- \\
Aegis~\cite{aegis}          & --          & --      & -- & -- & -- & \checkmark & -- \\
AgentDebug~\cite{agentdebug}& --          & --      & -- & -- & -- & -- & -- \\
\midrule
\multicolumn{8}{@{}l}{\emph{Debugging \& diagnosis tools}} \\
AgentRx~\cite{agentrx}      & $\circ$     & $\circ$ & \checkmark & -- & \checkmark & -- & -- \\
AGDebugger~\cite{agdebugger}& --          & $\circ$ & \checkmark & -- & $\circ$ & -- & -- \\
AgentDiagnose~\cite{agentdiagnose} & $\circ$ & -- & -- & -- & $\circ$ & \checkmark & -- \\
\midrule
\multicolumn{8}{@{}l}{\emph{Safety monitoring \& guardrails}} \\
GuardAgent~\cite{guardagent} & $\circ$    & --      & -- & -- & -- & \checkmark & -- \\
Runtime Enf.~\cite{runtime_enforcement} & -- & -- & -- & -- & -- & -- & -- \\
NeMo Guardrails~\cite{nemo_guardrails} & -- & --  & -- & -- & -- & \checkmark & -- \\
\midrule
\multicolumn{8}{@{}l}{\emph{LLM observability platforms}} \\
LangSmith~\cite{langsmith}  & --          & $\circ$ & -- & -- & -- & -- & -- \\
Langfuse~\cite{langfuse}    & $\circ$     & --      & -- & -- & -- & \checkmark & \checkmark \\
Datadog LLM~\cite{datadog_llm} & \checkmark & --    & -- & -- & -- & -- & $\circ$ \\
Weave~\cite{weave}          & $\circ$     & --      & -- & -- & -- & -- & -- \\
Arize Phoenix               & $\circ$     & --      & -- & -- & -- & \checkmark & \checkmark \\
Helicone                    & \checkmark  & --      & -- & -- & -- & -- & -- \\
OpenLLMetry~\cite{openllmetry} & \checkmark & --   & -- & -- & -- & \checkmark & \checkmark \\
\midrule
\multicolumn{8}{@{}l}{\emph{Agent SDKs \& protocols}} \\
OpenAI Agents~\cite{openai_agents_sdk} & -- & $\circ$ & $\circ$ & -- & -- & \checkmark & -- \\
MCP~\cite{mcp}              & \checkmark  & --      & -- & -- & -- & \checkmark & -- \\
\midrule
\multicolumn{8}{@{}l}{\emph{OTel standards}} \\
OTel GenAI~\cite{otel}      & \checkmark  & $\circ$ & $\circ$ & -- & -- & \checkmark & \checkmark \\
\midrule
\textbf{AgentTelemetry}     & \checkmark & \checkmark & \checkmark & \checkmark & \checkmark & \checkmark & \checkmark \\
\bottomrule
\end{tabular}
\end{table}

\subsection{Agent Failure Taxonomies}

Several recent works catalog how LLM agents fail.
MAST~\cite{mast} identifies 14 failure modes from 1{,}600+ annotated
multi-agent traces across 7~frameworks, organized into specification issues,
inter-agent misalignment, and task verification, with strong inter-annotator
agreement ($\kappa = 0.88$).
Table~\ref{tab:mast-mapping} maps MAST's failure modes to our fault types:
12 of~14 (85.7\%) are detectable via our span kinds, and the two gaps are
omission failures (failure to clarify, missing verification) that require
expected-span specifications beyond trace-level anomaly detection.
Dong et~al.~\cite{agentops} propose an observability taxonomy for agent
operations, identifying key telemetry dimensions but without an OTel-native
implementation or empirical fault-detection evaluation.
Aegis~\cite{aegis} classifies six agent-environment failure modes,
focusing on misalignment between agent actions and environmental
constraints.
AgentDebug~\cite{agentdebug} categorizes failure modes across four
cognitive dimensions---memory, reflection, planning, and action---with
a diagnostic framework for each.
These taxonomies characterize \emph{what} goes wrong but do not provide
runtime telemetry infrastructure to \emph{detect} faults in production.
\textsc{AgentTelemetry} bridges this gap: our 14 fault types
(Table~\ref{tab:faults}) are informed by these taxonomies, and our span
kinds provide the runtime observability primitives needed to detect them
automatically.

\subsection{Agent Debugging and Diagnosis Tools}

AgentRx~\cite{agentrx} is a Microsoft framework for automated
diagnosis of multi-agent systems, using LLM-based analysis of
execution logs to identify root causes.
AGDebugger~\cite{agdebugger} provides an interactive visual debugging
interface for multi-agent conversations, validated through a 14-person
user study that demonstrated improved debugging efficiency.
AgentDiagnose~\cite{agentdiagnose} offers a toolkit for systematic
analysis of agent trajectories, enabling developers to identify failure
patterns across execution traces.
These tools operate on post-hoc logs or require custom instrumentation;
\textsc{AgentTelemetry} provides the structured telemetry layer that
such tools could consume, offering standardized, OTel-native trace data
with agent-aware span kinds rather than unstructured logs.
The OpenAI Agents SDK~\cite{openai_agents_sdk} includes built-in
tracing with first-class concepts for handoffs and guardrails, but is
tied to a single vendor and does not export OTel-native spans or
define a portable span taxonomy.

\subsection{Agent Safety and Guardrails}

GuardAgent~\cite{guardagent} introduces an LLM-based guardrail agent
that dynamically generates and verifies safety constraints for other
agents, providing flexible runtime protection.
Runtime Enforcement for LLM Agents~\cite{runtime_enforcement} proposes
a domain-specific language for specifying and enforcing runtime safety
policies on agent behavior, validated across multiple agent benchmarks.
NeMo Guardrails~\cite{nemo_guardrails} enforces safety policies
through configurable dialog rails.

These systems are \emph{enforcement} mechanisms: they intercept agent
actions and block or modify unsafe behavior at runtime.
\textsc{AgentTelemetry}'s \texttt{GUARD\_RAIL} span kind serves a
fundamentally different role---it provides \emph{observability} for
safety checks, recording whether a guardrail fired, was bypassed, or
failed, but it does not itself enforce safety policies.  The
\texttt{GUARD\_RAIL} span is complementary to enforcement systems:
GuardAgent and NeMo Guardrails enforce policies, while
\texttt{GUARD\_RAIL} spans monitor enforcement outcomes.  This
distinction is critical because enforcement without monitoring creates
a blind spot: operators cannot verify that guardrails are functioning
correctly, detect bypass attempts, or audit policy compliance without
a dedicated observability primitive.  Conversely,
\texttt{GUARD\_RAIL} spans without an enforcement backend simply
record that no safety check occurred---useful for identifying gaps in
the safety infrastructure but not a substitute for enforcement.

\subsection{LLM Observability Platforms}

\textbf{Framework-specific platforms.}
LangSmith~\cite{langsmith} provides deep integration with LangChain
but no portability to other frameworks.
\textbf{Commercial and open-source platforms.}
Langfuse~\cite{langfuse} is an open-source LLM observability platform
with OTel integration for tracing, prompt management, and evaluation;
Datadog LLM Observability~\cite{datadog_llm} extends a commercial
monitoring stack with LLM-specific trace views and cost tracking;
Weave~\cite{weave} by Weights~\&~Biases provides production tracing
and evaluation for LLM applications.
These platforms capture LLM call traces and cost metrics but do not
define reusable semantic span kinds that compose with the broader OTel
ecosystem---they cannot distinguish a planning step from a
summarization call, and their trace formats are not portable across
backends.
\textbf{API-layer instrumentors.}
Helicone, OpenLLMetry~\cite{openllmetry}, and Arize Phoenix intercept
LLM API calls and provide token-level metrics, but cannot distinguish
agent-level semantics---a planning step and a summarization call are
both \texttt{POST /v1/chat/completions}.
\textbf{System-level monitoring.}
AgentSight~\cite{agentsight} uses eBPF to monitor TLS-encrypted LLM
traffic at the kernel level, providing deep visibility but no
application-level semantics (agent role, task decomposition, delegation
chains).
Zaharia et~al.~\cite{compound_ai} argue that compound AI systems
combining multiple LLM calls, retrievers, and tools require new
infrastructure beyond single-model serving; \textsc{AgentTelemetry}
provides this infrastructure for the observability dimension.

\subsection{OpenTelemetry Standards}

The OTel GenAI semantic conventions~\cite{otel} (all at Development
stability) define eight \texttt{gen\_ai.operation.name} values spanning
LLM invocation (\texttt{chat}, \texttt{text\_completion},
\texttt{generate\_content}), embedding and retrieval (\texttt{embeddings},
\texttt{retrieval}), tool execution (\texttt{execute\_tool}), and agent
lifecycle (\texttt{create\_agent}, \texttt{invoke\_agent}).  These
conventions also specify \texttt{gen\_ai.*} attributes for model, token
usage, and finish reason, plus \texttt{gen\_ai.agent.*} attributes for
agent name and ID.

\textsc{AgentTelemetry} is a \textbf{strict superset} of the GenAI
conventions for agent orchestration (Table~\ref{tab:otel_comparison}).
Three of our span kinds overlap with GenAI operations
(\texttt{LLM\_CALL}, \texttt{TOOL\_CALL}, \texttt{RETRIEVAL}); one
extends them (\texttt{AGENT} adds framework, role, and orchestration
context beyond \texttt{invoke\_agent}'s lifecycle tracking).  The
remaining \textbf{five span kinds are entirely novel}: \texttt{PLANNING},
\texttt{REASONING}, \texttt{GUARD\_RAIL}, \texttt{DELEGATION}, and
\texttt{MEMORY} have no GenAI equivalent---yet our ablation study
(\S\ref{sec:rq2}) proves each is necessary for detecting at least one
fault type.  These five kinds capture the agent orchestration
phases---task decomposition, chain-of-thought, safety-check monitoring,
inter-agent handoff, and state management---that distinguish autonomous
agents from simple LLM API wrappers.

\begin{table}[t]
\centering
\small
\caption{AgentTelemetry span kinds vs.\ OTel GenAI semantic conventions.
Three kinds overlap, one extends, and five are novel contributions
with no GenAI equivalent.  All five novel kinds are proven necessary
by our ablation study.}
\label{tab:otel_comparison}
\begin{tabular}{@{}lp{2.8cm}l@{}}
\toprule
\textbf{AgentTelemetry Kind} & \textbf{OTel GenAI Equivalent} & \textbf{Status} \\
\midrule
\texttt{LLM\_CALL}   & \texttt{chat}, \texttt{text\_completion}, \texttt{generate\_content} & Overlapping \\
\texttt{TOOL\_CALL}  & \texttt{execute\_tool}            & Overlapping \\
\texttt{RETRIEVAL}   & \texttt{retrieval}, \texttt{embeddings} & Overlapping \\
\texttt{AGENT}       & \texttt{create\_agent}, \texttt{invoke\_agent} & Extended \\
\midrule
\texttt{PLANNING}    & (none) & \textbf{Novel} \\
\texttt{REASONING}   & (none) & \textbf{Novel} \\
\texttt{GUARD\_RAIL} & (none) & \textbf{Novel} \\
\texttt{DELEGATION}  & (none) & \textbf{Novel} \\
\texttt{MEMORY}      & (none) & \textbf{Novel} \\
\bottomrule
\end{tabular}
\end{table}

\textsc{AgentTelemetry} is fully compatible with the GenAI conventions:
spans export via OTLP to any GenAI-aware backend, and GenAI attributes
(\texttt{gen\_ai.usage.*}, \texttt{gen\_ai.response.*}) can coexist with
our \texttt{agenttelemetry.*} attributes on the same spans.  Our design
thus provides an incremental adoption path: teams already using OTel
GenAI instrumentation gain the five novel orchestration span kinds
without abandoning existing telemetry pipelines.

Anthropic's Model Context Protocol (MCP)~\cite{mcp} standardizes how
AI assistants connect to external tools and data sources.  MCP defines
the \emph{tool interface} but not the \emph{observability layer}:
\textsc{AgentTelemetry}'s \texttt{TOOL\_CALL} and \texttt{DELEGATION}
spans can wrap MCP interactions, providing the trace-level visibility
that MCP itself does not specify.

%% ====================================================================
%% 8. CONCLUSION
%% ====================================================================
\section{Conclusion}
\label{sec:conclusion}

We have presented \textsc{AgentTelemetry}, the first OTel-native
observability library with a grounded taxonomy of agent-specific span
kinds that forms a strict superset of the OTel GenAI semantic
conventions for agent orchestration.  Five of nine span
kinds---\texttt{PLANNING}, \texttt{REASONING}, \texttt{GUARD\_RAIL},
\texttt{DELEGATION}, and \texttt{MEMORY}---are entirely novel, and
our ablation study proves that all nine are
necessary: removing any one makes at least one fault type undetectable.
The SWE-bench case study (RQ6) demonstrates that the taxonomy captures real failure modes on an established benchmark: the novel REASONING span kind characterizes the dominant failure mode (75\% of 112~instances, 95\% CI [66\%, 82\%]), which would be invisible under vanilla OTel or OTel GenAI conventions.  Moreover, a telemetry-guided intervention that detects and breaks reasoning loops improves the SWE-bench patch rate by $+11$~pp over a no-intervention control (3-seed mean), closing the loop from observability through diagnosis to measurable improvement in agent performance.
As a structural completeness check, our controlled fault-injection
evaluation across 14~fault types, 6~models, 7~frameworks, and
5~observability conditions confirms that the span taxonomy is
sufficient: FDR\,=\,1.000 for \textsc{AgentTelemetry} versus 0.429 for
vanilla OTel---an upper bound demonstrating that every fault type has a
corresponding detection signal.

Future work includes proposing our span taxonomy as an OpenTelemetry
Enhancement Proposal (OTEP), verifying SWE-bench patches against ground-truth test suites to measure resolve rate improvement, adaptive trace sampling for long-running
agents, real-time streaming analysis, and applying causal inference to
automated root cause analysis in multi-agent systems.

%% ====================================================================
%% BIBLIOGRAPHY
%% ====================================================================

\begin{thebibliography}{36}

\bibitem{agentsight}
Y.~Zheng, Y.~Hu, T.~Yu, and A.~Quinn.
\newblock {AgentSight}: System-Level Observability for {AI} Agents Using {eBPF}.
\newblock {\em arXiv preprint arXiv:2508.02736}, 2025.

\bibitem{nemo_guardrails}
T.~Rebedea, R.~Dinu, M.~Sreedhar, C.~Parisien, and J.~Cohen.
\newblock {NeMo Guardrails}: A Toolkit for Controllable and Safe {LLM}
  Applications.
\newblock In {\em Proc.\ EMNLP (Demo Track)}, 2023.

\bibitem{dspy_assertions}
A.~Singhvi, M.~Shetty, S.~Tan, C.~Potts, K.~Sen, M.~Zaharia, and O.~Khattab.
\newblock {DSPy} Assertions: Computational Constraints for Self-Refining
  Language Model Pipelines.
\newblock In {\em Proc.\ COLM}, 2024.

\bibitem{agentboard}
C.~Ma, J.~Zhang, Z.~Zhu, C.~Yang, Y.~Yang, et~al.
\newblock {AgentBoard}: An Analytical Evaluation Board of Multi-Turn {LLM}
  Agents.
\newblock In {\em Proc.\ ACL}, 2024.

\bibitem{agent_as_judge}
M.~Zhuge, C.~Zhao, D.~Ashley, et~al.
\newblock Agent-as-a-Judge: Evaluate Agents with Agents.
\newblock {\em arXiv preprint arXiv:2410.10934}, 2024.

\bibitem{agentbench}
X.~Liu, H.~Yu, H.~Zhang, et~al.
\newblock {AgentBench}: Evaluating {LLMs} as Agents.
\newblock In {\em Proc.\ ICLR}, 2024.

\bibitem{openagentsafety}
S.~Vijayvargiya, A.~B.~Soni, X.~Zhou, Z.~Z.~Wang, N.~Dziri, G.~Neubig, and
  M.~Sap.
\newblock {OpenAgentSafety}: A Comprehensive Framework for Evaluating
  Real-World {AI} Agent Safety.
\newblock {\em arXiv preprint arXiv:2507.06134}, 2025.

\bibitem{toolemu}
Y.~Ruan, H.~Dong, A.~Wang, S.~Pitis, et~al.
\newblock {ToolEmu}: Identifying Risks of {LLM} Agents with an {LLM}-Emulated
  Sandbox.
\newblock In {\em Proc.\ ICLR}, 2024.

\bibitem{ml_monitoring}
J.~Leest, I.~Gerostathopoulos, P.~Lago, and C.~Raibulet.
\newblock Monitoring and Observability of Machine Learning Systems: Current
  Practices and Gaps.
\newblock {\em arXiv preprint arXiv:2510.24142}, 2025.

\bibitem{aiops_survey}
L.~Zhang, T.~Jia, M.~Jia, Y.~Wu, A.~Liu, Y.~Yang, Z.~Wu, X.~Hu, P.~S.~Yu,
  and Y.~Li.
\newblock A Survey of {AIOps} in the Era of Large Language Models.
\newblock {\em arXiv preprint arXiv:2507.12472}, 2025.

\bibitem{otel}
{OpenTelemetry Authors}.
\newblock {OpenTelemetry Specification}, v1.30.
\newblock \url{https://opentelemetry.io/docs/specs/otel/}, 2024.

\bibitem{tracediag}
R.~Ding, C.~Zhang, L.~Wang, Y.~Xu, M.~Ma, W.~Zhang, S.~Qin, S.~Rajmohan,
  Q.~Lin, and D.~Zhang.
\newblock {TraceDiag}: Adaptive, Interpretable, and Efficient Root Cause
  Analysis on Large-Scale Microservice Systems.
\newblock In {\em Proc.\ ACM ESEC/FSE}, 2023.

\bibitem{deeptralog}
C.~Zhang, X.~Peng, C.~Sha, K.~Zhang, Z.~Fu, X.~Wu, Q.~Lin, and D.~Zhang.
\newblock {DeepTraLog}: Trace-Log Combined Microservice Anomaly Detection
  through Graph-based Deep Learning.
\newblock In {\em Proc.\ ACM ICSE}, 2022.

\bibitem{sieve_sampling}
D.~Batavia, I.~Weber, et~al.
\newblock {Sieve}: Attention-based Sampling of End-to-End Trace Data in
  Distributed Microservice Systems.
\newblock In {\em Proc.\ IEEE ICWS}, 2023.

\bibitem{mulan}
L.~Zheng, Z.~Chen, J.~He, and H.~Chen.
\newblock {MULAN}: Multi-modal Causal Structure Learning for Root Cause
  Analysis.
\newblock In {\em Proc.\ NeurIPS}, 2024.

\bibitem{langsmith}
{LangChain Inc.}
\newblock {LangSmith}: Platform for Building Production-Grade {LLM} Applications.
\newblock \url{https://smith.langchain.com/}, 2024.

\bibitem{openllmetry}
{Traceloop}.
\newblock {OpenLLMetry}: Open-Source Observability for {LLM} Applications
  Built on {OpenTelemetry}.
\newblock \url{https://github.com/traceloop/openllmetry}, 2024.

\bibitem{strauss_corbin}
A.~Strauss and J.~Corbin.
\newblock {\em Basics of Qualitative Research: Techniques and Procedures for
  Developing Grounded Theory}.
\newblock Sage, 2nd edition, 1998.

\bibitem{agentdebug}
M.~Cemri, M.~Y.~Qiang, Y.~Xiao, and L.~Tan.
\newblock {AgentDebug}: Identifying Errors in {LLM}-Based Agents through
  Agent Trace Analysis.
\newblock {\em arXiv preprint arXiv:2504.16744}, 2025.

\bibitem{mast}
M.~Cemri, M.~Z.~Pan, S.~Yang, L.~A.~Agrawal, B.~Chopra, R.~Tiwari,
  K.~Keutzer, A.~Parameswaran, D.~Klein, K.~Ramchandran, M.~Zaharia,
  J.~E.~Gonzalez, and I.~Stoica.
\newblock Why Do Multi-Agent {LLM} Systems Fail?
\newblock In {\em Proc.\ NeurIPS}, 2025.

\bibitem{agentops}
J.~Dong, Y.~Liu, H.~Qian, L.~Zheng, and L.~Chen.
\newblock {AgentOps}: An Observability Taxonomy for {AI} Agent Operations.
\newblock {\em arXiv preprint arXiv:2411.05285}, 2024.

\bibitem{agentrx}
{Microsoft Research}.
\newblock {AgentRx}: Automated Diagnostic Framework for Multi-Agent Systems.
\newblock {\em arXiv preprint arXiv:2602.02475}, 2026.

\bibitem{agdebugger}
S.~Jiang, X.~Wang, Y.~Zhang, et~al.
\newblock {AGDebugger}: An Interactive Multi-Agent Debugging Tool.
\newblock In {\em Proc.\ ACM CHI}, 2025.

\bibitem{agentdiagnose}
Y.~Li, Z.~Wang, H.~Chen, et~al.
\newblock {AgentDiagnose}: A Toolkit for Diagnosing Agent Trajectories.
\newblock In {\em Proc.\ EMNLP (Demo Track)}, 2025.

\bibitem{aegis}
Z.~Chen, L.~Xu, M.~Wang, et~al.
\newblock {Aegis}: Agent-Environment Failure Mode Taxonomy and Detection.
\newblock {\em arXiv preprint arXiv:2508.19504}, 2025.

\bibitem{guardagent}
X.~Zhan, Y.~Luo, S.~Wang, et~al.
\newblock {GuardAgent}: Safeguard {LLM} Agents by a Guard Agent via
  Knowledge-Enabled Reasoning.
\newblock In {\em Proc.\ ICML}, 2025.

\bibitem{runtime_enforcement}
A.~Roy, S.~Elnikety, A.~Sarkar, and Y.~Smaragdakis.
\newblock Runtime Enforcement for {LLM} Agents.
\newblock In {\em Proc.\ ACM ICSE}, 2026.

\bibitem{compound_ai}
M.~Zaharia, O.~Khattab, L.~Chen, J.~Q.~Davis, H.~Ge, et~al.
\newblock The Shift from Models to Compound {AI} Systems.
\newblock Berkeley AI Research Blog, 2024.

\bibitem{otel_agent_spans}
{OpenTelemetry Authors}.
\newblock {GenAI} Agent Spans Semantic Conventions, v1.37.0+.
\newblock \url{https://opentelemetry.io/docs/specs/semconv/gen-ai/gen-ai-agent-spans/}, 2025.

\bibitem{langfuse}
{Langfuse GmbH}.
\newblock {Langfuse}: Open-Source {LLM} Observability Platform.
\newblock \url{https://langfuse.com/}, 2024.

\bibitem{datadog_llm}
{Datadog Inc.}
\newblock {LLM Observability}: Monitor and Improve {LLM} Applications.
\newblock \url{https://www.datadoghq.com/product/llm-observability/}, 2024.

\bibitem{weave}
{Weights \& Biases}.
\newblock {Weave}: Production Tracing and Evaluation for {LLM} Applications.
\newblock \url{https://wandb.ai/site/weave/}, 2024.

\bibitem{openai_agents_sdk}
{OpenAI}.
\newblock {OpenAI Agents SDK}: Build Multi-Agent Workflows with Built-in
  Tracing.
\newblock \url{https://github.com/openai/openai-agents-python}, 2025.

\bibitem{mcp}
{Anthropic}.
\newblock {Model Context Protocol ({MCP})}: An Open Standard for Connecting
  {AI} Assistants to External Tools and Data.
\newblock \url{https://modelcontextprotocol.io/}, 2024.

\bibitem{swebench}
C.~E. Jimenez, J.~Yang, A.~Wettig, S.~Yao, K.~Pei, O.~Press, and
K.~Narasimhan.
\newblock {SWE-bench}: Can Language Models Resolve Real-World {GitHub} Issues?
\newblock In \emph{ICLR}, 2024.

\end{thebibliography}

\end{document}
